%6622770764
\documentclass{ieeetj}

% ---------------------------------------------------------
% Packages
% ---------------------------------------------------------
\usepackage{placeins}
\usepackage{cite}
\usepackage{amsmath,amssymb,amsfonts}
\usepackage{amsthm}
\usepackage{algorithm}
\usepackage{algpseudocode}
\usepackage{graphi  cx}
\usepackage{xcolor}
\usepackage{textcomp}
\usepackage{hyperref}
\hypersetup{hidelinks=true}
\usepackage{booktabs}
\usepackage{tabularx}
\usepackage{array}
\usepackage{makecell}
\usepackage{multirow}
\usepackage{amssymb}
\usepackage{pifont}
\newcommand{\cmark}{\ding{51}}
\newcommand{\xmark}{\ding{55}}
\newcommand{\pmark}{\ding{108}}

% ---------------------------------------------------------
% BibTeX macro (correct)
% ---------------------------------------------------------
\def\BibTeX{{\rm B\kern-.05em{\sc i\kern-.025em b}\kern-.08em
 T\kern-.1667em\lower.7ex\hbox{E}\kern-.125emX}}

% ---------------------------------------------------------
% Theorem Style
% ---------------------------------------------------------
\newtheoremstyle{ieee-style}%
  {\topsep}{\topsep}%
  {\itshape}{}%
  {\bfseries\itshape}{.}{ }{}%

\theoremstyle{ieee-style}

% ---------------------------------------------------------
% Theorem Environments
% ---------------------------------------------------------
\newtheorem{theorem}{Theorem}[subsection]
\newtheorem{lemma}{Lemma}[subsection]
\newtheorem{definition}{Definition}[subsection]
\newtheorem{corollary}{Corollary}[subsection]

% QED symbol
\renewcommand{\qedsymbol}{$\square$}

% ---------------------------------------------------------
% Color Definitions
% ---------------------------------------------------------
\AtBeginDocument{
  \definecolor{tmlcncolor}{cmyk}{0.93,0.59,0.15,0.02}
  \definecolor{NavyBlue}{RGB}{0,86,125}
}





\def\OJlogo{\vspace{-4pt}$<$Society logo(s) and publication title will appear here.$>$}
\def\seclogo{\vspace{10pt}$<$Society logo(s) and publication title will appear here.$>$}

\def\authorrefmark#1{\ensuremath{^{\textbf{#1}}}}

\begin{document}
\receiveddate{XX Month, XXXX}
\reviseddate{XX Month, XXXX}
\accepteddate{XX Month, XXXX}
\publisheddate{XX Month, XXXX}
\currentdate{XX Month, XXXX}
\doiinfo{XXXX.2022.1234567}

\markboth{}{Fugkeaw \MakeLowercase{\textit{et al.}}}

\title{{Achieving Post-Quantum and Dynamic Load-Balanced Verifiable Searchable Encryption for Multi-Authority IoMT Data Sharing}}

\author{
Somchart Fugkeaw\authorrefmark{1},
Author2,Author3\authorrefmark{1},\\
}


\affil{Sirindhorn International Institute of Technology, Thammasat University, Pathum Thani, Thailand 

somchart@siit.tu.ac.th,
pattavee.san@gmail.com}


\corresp{Corresponding author:  Somchart Fugkeaw (somchart@siit.tu.ac.th).}


\begin{abstract}
Large-scale Internet of Medical Things (IoMT) systems continuously generate sensitive healthcare data that must be securely shared across multiple administrative domains. Existing searchable encryption schemes largely assume centralized search infrastructures and static authorization models, resulting in inefficient encrypted retrieval, costly synchronization, and limited scalability under dynamic IoMT workloads. Moreover, they rarely integrate multi-authority authorization, verifiable retrieval, and adaptive search scheduling within a unified framework. This paper proposes \textbf{MA-LB-PQ-VDSE} (\emph{Multi-Authority Load-Balanced Post-Quantum Verifiable Dynamic Searchable Encryption}), a secure and scalable framework for dynamic IoMT data sharing. The proposed framework unifies multi-authority authorization, policy-bound dynamic searchable encryption, adaptive authorization-aware search scheduling, incremental authorization synchronization, blockchain-backed verifiable retrieval, and post-quantum secure key establishment. By tightly coupling authorization management, encrypted search, dynamic updates, and verification, MA-LB-PQ-VDSE significantly reduces unnecessary index traversal, synchronization overhead, and encrypted-search latency while preserving authorization consistency and retrieval integrity. Furthermore, an adaptive scheduler predicts cryptographic search costs to efficiently distribute encrypted queries across Fog Search Nodes, improving scalability under bursty workloads. Experimental results show that the proposed framework consistently outperforms existing schemes in search latency, throughput, synchronization efficiency, and scalability, demonstrating its practicality for large-scale cross-domain IoMT data sharing.
\end{abstract}


\begin{keywords}
EHRs, dynamic symmetric searchable encryption, bloom filter, post-quantum, blockchain, foreward privacy, backword privacy.
\end{keywords}



\maketitle

\section{Introduction}
\label{sec:introduction}

\PARstart{T}{he} Internet of Medical Things (IoMT) enables continuous
collection and sharing of medical data from wearable devices, clinical
sensors, imaging systems, ambulances, and hospital information systems.
Such data are increasingly outsourced to cloud--fog infrastructures to
support telemedicine, collaborative diagnosis, emergency care, and
large-scale healthcare analytics. However, outsourcing sensitive medical
records to semi-trusted infrastructures introduces significant challenges
in confidentiality, fine-grained cross-domain authorization, secure
retrieval, integrity verification, and regulatory compliance with HIPAA,
GDPR, and PDPA~\cite{ref1,ref2,ref3,ref4,ref5,ref20,ref21,ref47,ref50,ref51}.

Searchable encryption enables users to query encrypted outsourced data
without first decrypting the entire repository. Dynamic Searchable
Symmetric Encryption (DSSE) further supports record insertion, deletion,
and modification~\cite{ref29,ref30,ref31,ref37,ref38,ref42,ref45,ref46},
while verifiable searchable encryption employs authenticated structures
such as Merkle trees and blockchain commitments to detect forged or
modified search evidence returned by untrusted servers
~\cite{ref6,ref8,ref9,ref10,ref22,ref23,ref24,ref30,ref31,ref32,ref35,ref46}.
In parallel, multi-authority attribute-based mechanisms distribute
authorization across independent administrative domains
~\cite{ref26,ref27,ref28}, and recent LWE- and lattice-based approaches
investigate post-quantum protection for searchable data sharing
~\cite{ref16,ref41,ref52,ref54,ref55}. These advances address important
cryptographic requirements, but they do not fully resolve how
authorization evolution, dynamic encrypted indexes, distributed search,
and verifiable retrieval should interact in a large multi-domain IoMT
deployment.

A fundamental challenge is that \emph{authorization and searchable-index
states evolve together but at different scopes}. Consider encrypted IoMT
records distributed across hospital, laboratory, and emergency-care
domains. A cardiologist may query \emph{``ECG records associated with
arrhythmia''}, while an emergency physician may search \emph{``cardiac
records requiring emergency intervention''} across multiple domains.
Such queries should access only shards whose policies are satisfied and
whose relevant authority states are current. However, state-independent
tokens may permit stale authorization, whereas binding indexes to a
global authorization version causes an unrelated AA update to invalidate
unaffected indexes. This motivates policy-state-aware tokens and
commitments bound only to the AAs governing each policy.

A second challenge is efficient distributed encrypted search. A
cross-domain query may span shards maintained by different Fog Search
Nodes (FSNs), while the least-loaded FSN may lack a required shard,
maintain stale authority state, or incur high index and proof-generation
costs. Resource-only load balancing therefore poorly reflects encrypted
search overhead. Scheduling should instead consider \emph{index locality,
policy-state synchronization, verification workload, and queue delay}
when assigning authorized shards.

IoMT data and authorization are also inherently dynamic: records are
inserted, modified, or deleted, while sharing policies and user attributes
may change or be revoked. These events have different scopes. A record
modification changes its encrypted content and searchable entries, whereas
an attribute revocation may leave ciphertexts unchanged but invalidate
indexes governed by the affected AA. For example, revocation by a research
AA should not update emergency records governed only by hospital and
emergency AAs. Global re-keying, index reconstruction, or FSN-wide
synchronization therefore introduces unnecessary overhead.

Finally, dynamic evolution complicates verifiable retrieval. A record
modification produces a new content-addressed encrypted object, whereas a
policy or revocation change may require only updated policy-state tokens,
index roots, and commitments. A returned $CID_i$ alone cannot prove that
its index entry belongs to the current authenticated index or policy
state. Hence, the framework must distinguish \emph{data evolution} from
\emph{authorization evolution}, update only affected dependencies, and
cryptographically bind searchable state to the exact encrypted object and
current policy state.
To address these challenges, we propose
\textbf{MA-LB-PQ-VDSE}
(\emph{Multi-Authority Load-Balanced Post-Quantum Verifiable Dynamic
Searchable Encryption}), a cloud--fog framework for scalable multi-domain
IoMT retrieval. MA-LB-PQ-VDSE decouples content-addressed encrypted
records from a policy-state-aware dynamic index. A Version-Bound
Authorization Profile (VAP) captures user-specific multi-authority state,
while policy-state digests bind search tokens only to the AAs governing
each policy, preventing unrelated AA updates from invalidating unaffected
indexes. The Adaptive Authorization-Aware Search Scheduler (AASS) assigns
authorized shards according to cryptographic, synchronization, and queue
costs, while Dependency-Aware Incremental Authorization Synchronization
(DIAS) updates only affected policies, records, shards, users, and FSNs.
Merkle proofs and blockchain commitments bind returned index entries to
their encrypted objects and current policy states, while ML-KEM provides
post-quantum protection for attribute-key delivery.

The main contributions are summarized as follows:

\begin{itemize}

\item \textbf{Policy-state-aware multi-authority dynamic search.}
We develop VAP and a policy-state-aware dynamic index that bind search
tokens and authenticated commitments only to the current states of
policy-governing AAs. This supports authorized cross-domain search while
preventing unrelated authority updates from invalidating unaffected index
entries.

\item \textbf{Adaptive authorization-aware distributed scheduling.}
We design AASS to assign each required index shard to an eligible Fog
Search Node according to predicted index traversal, proof generation,
policy-state synchronization, and queue costs, thereby incorporating
cryptographic and authorization overhead beyond resource-only load
balancing.

\item \textbf{Dependency-aware dynamic evolution and synchronization.}
We introduce DIAS to support data and authorization changes according to
their dependency scope. It incrementally updates only affected policies,
records, index shards, users, and FSNs, avoiding global index
reconstruction and authorization-state redistribution.

\item \textbf{Decoupled verifiable outsourcing and protected key delivery.}
We decouple encrypted IoMT objects from evolvable search indexes using
content-addressed references, allowing policy changes without
re-outsourcing unchanged ciphertexts. ML-KEM protects attribute-key
delivery, while authenticated encryption, Merkle proofs, and
blockchain-anchored commitments provide integrity- and
authorization-consistent retrieval.

\end{itemize}


\section{Related Work}
\begin{table*}[t]
\centering
\caption{Comparison of representative searchable encryption frameworks.}
\label{tab:comparison}
\renewcommand{\arraystretch}{1.1}
\setlength{\tabcolsep}{5pt}

\begin{tabular}{lccccccc}
\hline
\textbf{Reference} &
\textbf{Dynamic} &
\textbf{Verifiable} &
\textbf{Multi-} &
\textbf{Blockchain} &
\textbf{Multi-Keyword} &
\textbf{Load} &
\textbf{Lattice/PQ} \\
&
\textbf{Update} &
\textbf{Search} &
\textbf{Authority} &
\textbf{Audit} &
\textbf{Search} &
\textbf{Balancing} &
\textbf{Based} \\
\hline

LWE-based MWSE~\cite{ref16}
& $\times$ & $\times$ & $\times$ & $\times$ & $\times$ & $\times$ & $\checkmark$ \\

Xu \textit{et al}.~\cite{ref26}
& $\times$ & $\times$ & $\checkmark$ & $\times$ & $\triangle$ & $\times$ & $\times$ \\

Yu \textit{et al}.~\cite{ref28}
& $\triangle$ & $\times$ & $\checkmark$ & $\checkmark$ & $\times$ & $\times$ & $\times$ \\

Dou \textit{et al}.~\cite{ref29}
& $\checkmark$ & $\times$ & $\times$ & $\times$ & $\times$ & $\times$ & $\triangle$ \\

Ge \textit{et al}.~\cite{ref30}
& $\checkmark$ & $\checkmark$ & $\times$ & $\checkmark$ & $\times$ & $\times$ & $\times$ \\

Guo \textit{et al}.~\cite{ref35}
& $\checkmark$ & $\checkmark$ & $\times$ & $\times$ & $\checkmark$ & $\times$ & $\times$ \\

Thingom \textit{et al.}~\cite{ref41}
& $\times$ & $\times$ & $\times$ & $\times$ & $\triangle$ & $\times$ & $\triangle$ \\

Chen \textit{et al}.~\cite{ref46}
& $\checkmark$ & $\checkmark$ & $\times$ & $\times$ & $\checkmark$ & $\times$ & $\times$ \\

Gao \textit{et al}.~\cite{ref47}
& $\times$ & $\triangle$ & $\times$ & $\checkmark$ & $\times$ & $\times$ & $\times$ \\

Perera and Fugkeaw~\cite{ref48}
& $\triangle$ & $\triangle$ & $\times$ & $\checkmark$ & $\checkmark$ & $\checkmark$ & $\times$ \\

Deebak and Hwang~\cite{ref51}
& $\times$ & $\times$ & $\times$ & $\checkmark$ & $\times$ & $\times$ & $\times$ \\

Zhuang \textit{et al.}~\cite{ref52}
& $\triangle$ & $\times$ & $\checkmark$ & $\times$ & $\checkmark$ & $\times$ & $\checkmark$ \\

Chen \textit{et al.}~\cite{ref53}
& $\times$ & $\times$ & $\checkmark$ & $\times$ & $\triangle$ & $\times$ & $\times$ \\

Perera and Fugkeaw~\cite{ref54}
& $\triangle$& $\checkmark$ & $\times$ & $\checkmark$ & $\checkmark$ & $\times$ & $\checkmark$ \\
\hline
\textbf{Proposed}
& $\checkmark$
& $\checkmark$
& $\checkmark$
& $\checkmark$
& $\checkmark$
& $\checkmark$
& $\checkmark$ \\
\hline
\end{tabular}
\end{table*}


\subsection{Blockchain-Assisted Secure Healthcare and IoMT Data Sharing}

Blockchain has been widely adopted to provide decentralized trust, immutable auditing, and fine-grained access control for healthcare and IoMT data sharing. Malamas \textit{et al.}~\cite{ref1}, Gao \textit{et al.}~\cite{ref2}, Chen \textit{et al.}~\cite{ref4}, and Yuan \textit{et al.}~\cite{ref5} proposed blockchain-based architectures for secure medical data sharing and decentralized access management. Fugkeaw~\cite{ref3} introduced an efficient policy update mechanism for outsourced health records, while Fugkeaw \textit{et al.}~\cite{ref21} incorporated adaptive revocation and fog-assisted load sharing for IoT healthcare systems. More recently, Gao \textit{et al.}~\cite{ref47}, Hak and Fugkeaw~\cite{ref50}, and Deebak and Hwang~\cite{ref51} extended blockchain-assisted healthcare sharing toward cloud-edge IoMT environments. Although these approaches improve decentralized trust and access control, they primarily focus on secure storage and authorization, offering limited support for efficient dynamic searchable encryption and scalable encrypted retrieval across multiple administrative domains.

\subsection{Dynamic and Verifiable Searchable Encryption}

Searchable Symmetric Encryption (SSE) originated with Song \textit{et al.}~\cite{ref39} and was formally strengthened by Curtmola \textit{et al.}~\cite{ref40}, while Thingom~\cite{ref41} introduced forward-secure searchable encryption. Recent Dynamic Searchable Symmetric Encryption (DSSE) schemes have significantly improved dynamic updates, forward/backward privacy, expressive queries, and multi-user support~\cite{ref29,ref33,ref35,ref36,ref37,ref38,ref42,ref45}. To ensure search correctness, numerous verifiable searchable encryption (VSE) schemes have integrated Merkle authentication, authenticated data structures, blockchain auditing, and cryptographic proofs~\cite{ref7,ref8,ref9,ref10,ref11,ref12,ref13,ref14,ref22,ref23,ref24,ref25,ref30,ref31,ref32,ref43,ref46}. These studies substantially improve security and verifiability; however, search execution generally remains centralized, verification is largely decoupled from authorization management, and the schemes do not consider workload-aware encrypted search in large-scale IoMT environments.

\subsection{Multi-Authority Authorization and Searchable Encryption}

Fine-grained authorization across multiple organizations has been extensively studied using Attribute-Based Encryption (ABE). Bethencourt \textit{et al.}~\cite{ref27} first introduced CP-ABE, which has become the foundation of secure healthcare data sharing. Xu \textit{et al.}~\cite{ref26} proposed authorized encrypted search for multi-authority medical databases, while Yu \textit{et al.}~\cite{ref28} combined blockchain with revocable multi-authority ABE for Industrial IoT. Other studies explored certificateless searchable encryption~\cite{ref19}, LWE-based multi-writer searchable encryption~\cite{ref16}, multi-user searchable encryption~\cite{ref17}, divertible searchable encryption~\cite{ref18}, and multi-client encrypted search~\cite{ref15}. Despite these advances, existing multi-authority schemes primarily concentrate on cryptographic authorization and key management, while overlooking dynamic searchable encryption over continuously generated IoMT data and intelligent workload balancing across distributed fog-cloud infrastructures.

Overall, existing research has advanced blockchain-assisted healthcare sharing, dynamic and verifiable searchable encryption, and multi-authority authorization independently. However, no existing framework jointly integrates multi-authority authorization, dynamic searchable encryption, verifiable encrypted retrieval, blockchain-backed auditing, and adaptive load-balanced search execution for large-scale IoMT data sharing. These limitations motivate the proposed framework.

Table~\ref{tab:comparison} compares representative searchable encryption frameworks with the proposed scheme. Existing studies have addressed dynamic searchable encryption, verifiable retrieval, multi-authority authorization, blockchain-assisted auditing, or lattice-/post-quantum cryptography individually, but none integrates these capabilities within a unified framework. Xu \textit{et al}.~\cite{ref26} and Yu \textit{et al}.~\cite{ref28} support multi-authority access control but lack verifiable encrypted search and workload-aware search execution. Representative DSSE and VDSE schemes~\cite{ref29,ref35,ref46} improve dynamic updates, search privacy, and retrieval verification, yet generally assume centralized search processing and do not provide blockchain-backed auditing or quantum-resistant security. Recent lattice-/post-quantum approaches~\cite{ref16,ref41,ref52,ref54} strengthen cryptographic resilience against quantum adversaries but primarily focus on secure searchable-encryption constructions rather than scalable authorization management, verifiable retrieval, and distributed search orchestration. Blockchain-assisted healthcare frameworks~\cite{ref47,ref51} enhance decentralized trust and auditability but do not jointly support dynamic searchable encryption and verifiable retrieval, while the load-balanced framework of Perera and Fugkeaw~\cite{ref48} improves search scalability without supporting multi-authority authorization, verifiable retrieval, or quantum-resistant cryptography. In contrast, the proposed framework is the first to unify dynamic searchable encryption, verifiable retrieval, multi-authority authorization, blockchain-backed auditing, lattice-/post-quantum secure communication, expressive Boolean search, and adaptive authorization-aware load-balanced search execution within a single architecture for large-scale cross-domain IoMT data sharing.
\section{Our Proposed System}

This section presents the architecture of
\textbf{MA-LB-PQ-VDSE}, a multi-authority,
load-balanced, post-quantum verifiable dynamic searchable
encryption framework for large-scale IoMT data sharing.
The framework enables fine-grained authorization,
dynamic encrypted search, efficient index maintenance,
verifiable retrieval, and adaptive workload distribution
across distributed fog-cloud infrastructures.

\subsection{System Model}

Fig.~\ref{fig:system} illustrates the MA-LB-PQ-VDSE architecture for
secure and verifiable IoMT data sharing across multiple healthcare
domains. The framework integrates multi-authority authorization,
policy-state-aware encrypted search, adaptive fog scheduling, dynamic
synchronization, and blockchain-assisted verification.

\begin{figure*}[!t]
    \centering
    \includegraphics[width=\textwidth]{MA-LB-PQ-VDSE System Model.png}
    \caption{Overall architecture of the proposed MA-LB-PQ-VDSE framework.}
    \label{fig:system}
\end{figure*}

\begin{itemize}

\item \textbf{IoMT Devices and Gateways:}
IoMT devices generate healthcare data, while edge gateways aggregate and
preprocess the collected records before secure outsourcing.

\item \textbf{Data Owner (DO):}
A healthcare provider acting as the DO encrypts IoMT records, defines
multi-authority access policies, outsources encrypted objects to IPFS,
and constructs policy-state-aware searchable indexes and commitments.

\item \textbf{Data User (DU):}
Authorized healthcare personnel generate search requests, verify returned
search evidence, retrieve encrypted records, and decrypt data when their
multi-authority attributes satisfy the corresponding access policies.

\item \textbf{Attribute Authorities (AAs):}
Independent AAs manage domain-specific attributes, user registration,
attribute-key issuance, revocation, and authorization-state commitments.
Their independent states collectively support cross-domain authorization
without a centralized authority.

\item \textbf{Authorization Index Manager (AIM):}
The AIM maintains user authorization metadata, policy mappings, VAPs, and
current multi-authority states. It validates search authorization, derives
policy-state-aware query tokens and authorized shard sets, and coordinates
dependency-aware state updates through DIAS.

\item \textbf{Load Balancing Controller (LBC):}
The LBC executes AASS to assign required search shards to eligible FSNs
according to index locality, policy-state synchronization, predicted
verification workload, and queue delay.

\item \textbf{Fog Search Nodes (FSNs):}
Distributed FSNs maintain searchable-index shards, execute encrypted
search over assigned authorized shards, synchronize affected states, and
generate Merkle authentication evidence for returned index entries.

\item \textbf{InterPlanetary File System (IPFS):}
IPFS stores encrypted IoMT objects off-chain using content-addressed
identifiers ($CID$s), decoupling large ciphertext storage from the
evolvable searchable index and blockchain metadata.

\item \textbf{Consortium Blockchain:}
The blockchain anchors authority states, policy-state-aware index
commitments, Merkle roots, and dynamic-update metadata. It provides an
authenticated reference for authorization-state freshness and
search-result verification without storing IoMT ciphertexts on-chain.

\end{itemize}

\section{MA-LB-PQ-VDSE Protocol}

This section presents the proposed MA-LB-PQ-VDSE protocol for secure and scalable IoMT data sharing across multiple administrative domains. Table~\ref{tab:notation} summarizes the primary notations used throughout the protocol.

\begin{table}[t]
\caption{Summary of Major Notations}
\label{tab:notation}
\centering
\footnotesize
\renewcommand{\arraystretch}{1.05}
\setlength{\tabcolsep}{3pt}

\begin{tabular}{p{0.30\columnwidth}p{0.56\columnwidth}}
\toprule
\textbf{Symbol} & \textbf{Description} \\
\midrule

$AA_i$
& $i$th Attribute Authority \\

$PK_i,MSK_i$
& Public and master secret keys of $AA_i$ \\

$SK_u$
& User attribute secret key \\

$M_i$
& IoMT data record \\

$CT_i$
& Encrypted IoMT ciphertext \\

$I_i$
& Dynamic searchable index \\

$CID_i$
& IPFS content identifier \\

$PID_i$
& Access-policy identifier \\

$VID_i$
& Version identifier \\

$ST$
& Search token \\

$Root_t$
& Merkle root at epoch $t$ \\

$\Pi$
& Verification proof \\

$AIM$
& Authorization Index Manager \\

$LBC$
& Load Balancing Controller \\

$FSN_j$
& $j$th Fog Search Node \\

$Score_j$
& Load-balancing score of $FSN_j$ \\

\bottomrule
\end{tabular}
\end{table}


\subsubsection*{\textbf{Phase I: System Initialization}}

This phase establishes the global cryptographic and system parameters for
multi-authority authorization, dynamic searchable encryption,
post-quantum key delivery, authenticated data encryption, verifiable
retrieval, and blockchain-assisted auditing. A trusted initializer
configures the common parameters, while each Attribute Authority (AA)
independently generates its authorization credentials. The initializer
is required only during system bootstrapping and does not participate in
subsequent authorization or search operations.

\paragraph*{Step 1: Global Cryptographic Initialization}

Given security parameter $\lambda$, the trusted initializer selects
bilinear groups $G_1$, $G_2$, and $G_T$ of prime order $p$, a
non-degenerate bilinear map
\begin{equation}
e:G_1\times G_2\rightarrow G_T
\end{equation}
and generators $g_1\in G_1$ and $g_2\in G_2$ for the underlying
multi-authority attribute-based authorization mechanism.

The cryptographic primitive set is defined as
\begin{equation}
\mathcal{P}=
\left(
H,\mathsf{AE},\mathsf{KDF},\mathsf{KEM},\mathsf{Merkle}
\right)
\end{equation}
where $H:\{0,1\}^{*}\rightarrow\{0,1\}^{\lambda}$ is a
collision-resistant hash function; $\mathsf{AE}$ denotes AES-256-GCM;
$\mathsf{KDF}$ denotes HKDF; $\mathsf{KEM}$ denotes ML-KEM for
post-quantum protected key delivery; and $\mathsf{Merkle}$ denotes the
hash-based authenticated-tree construction used for verifiable index
commitments. In implementation, $H$ may be instantiated using SHA-256.

For all structured objects, $\mathsf{Encode}(\cdot)$ denotes a
deterministic canonical encoding before hashing, ensuring consistent
commitment and verification across system entities.

\paragraph*{Step 2: Multi-Authority Initialization}

Let
\begin{equation}
\mathcal{AA}
=
\{AA_1,AA_2,\ldots,AA_{N_{AA}}\}
\end{equation}
denote the participating AAs. Each $AA_i$ independently executes
\begin{equation}
(MSK_i,PK_i)
\leftarrow
\mathsf{AA.Setup}(1^\lambda)
\end{equation}
where
\begin{equation}
MSK_i=(\alpha_i,\beta_i)
\end{equation}
and
\begin{equation}
PK_i=
\left(
g_1,g_2,
e(g_1,g_2)^{\alpha_i},
g_1^{\beta_i}
\right)
\end{equation}
with $\alpha_i,\beta_i\overset{\$}{\leftarrow}\mathbb{Z}_p$.
Each authority retains $MSK_i$ locally and publishes only $PK_i$.
Thus, no AA possesses another authority's master secret, enabling
independent cross-domain attribute management without a centralized
authorization authority.

\paragraph*{Step 3: Key-Delivery and Search-Infrastructure Initialization}

The system enables each enrolling user $U$ to generate an ML-KEM key pair
\begin{equation}
(pk_U^{\mathsf{KEM}},sk_U^{\mathsf{KEM}})
\leftarrow
\mathsf{KEM.KeyGen}(1^\lambda)
\end{equation}
which is subsequently used in Phase~III to protect the delivery of
attribute-key material.

The cloud--fog search infrastructure comprises the Authorization Index
Manager (AIM), Load Balancing Controller (LBC), and Fog Search Nodes
\begin{equation}
\mathcal{F}
=
\{FSN_1,FSN_2,\ldots,FSN_m\}
\end{equation}
Each $FSN_j$ maintains
\begin{equation}
\mathsf{State}_{FSN_j}
=
(\mathcal{S}_j,\mathbf{V}_j,Q_j)
\end{equation}
where $\mathcal{S}_j$ denotes its searchable-index shards,
$\mathbf{V}_j$ its latest synchronized multi-authority version state,
and $Q_j$ its current search-queue state. These states support subsequent
authorization-aware scheduling and incremental synchronization.

\paragraph*{Step 4: Global Public Parameter Publication}

The global public parameters are
\begin{equation}
\mathcal{PP}=
\left(
\lambda,
G_1,G_2,G_T,e,g_1,g_2,
\{PK_i\}_{i=1}^{N_{AA}},
\mathcal{P},
N_{AA},m
\right)
\end{equation}
The initializer publishes $\mathcal{PP}$ to all authorized system
entities. The consortium blockchain initializes an authenticated system
registry for subsequently anchoring authority registrations,
authorization-state commitments, policy-state-aware index commitments,
dynamic-update states, and audit records.

\subsubsection*{\textbf{Phase II: Multi-Authority Registration and
Authorization-State Commitment}}

Each Attribute Authority (AA) registers its public credentials and
establishes a blockchain-anchored commitment to its attribute namespace,
authorization version, and revocation state. These commitments form the
authenticated multi-authority state used for subsequent authorization,
policy-state-aware indexing, search, and synchronization.

\paragraph*{Step 1: Authority Registration}

Each $AA_i\in\mathcal{AA}$ publishes
\begin{equation}
Reg_i=(ID_i,Dom_i,PK_i)
\end{equation}
where $ID_i$, $Dom_i$, and $PK_i$ denote its unique identifier,
administrative domain, and public key generated in Phase~I, respectively.
The consortium blockchain anchors $Reg_i$, allowing system entities to
authenticate the issuing authority.

\paragraph*{Step 2: Attribute Namespace Initialization}

Each $AA_i$ independently maintains
\begin{equation}
\mathcal{A}_i=\{a_{i,1},a_{i,2},\ldots,a_{i,n_i}\}
\end{equation}
where attributes are uniquely identified by $(ID_i,a_{i,j})$ to avoid
cross-domain naming ambiguity. The authority computes
\begin{equation}
AHash_i=
H\!\left(\mathsf{Encode}(\mathcal{A}_i)\right)
\end{equation}
where $\mathsf{Encode}(\cdot)$ is the canonical encoding defined in
Phase~I.

\paragraph*{Step 3: Authorization-State Commitment}

Initially, $AA_i$ sets $v_i^{(0)}=0$ and constructs
\begin{equation}
RevRoot_i^{(0)}
=
\mathsf{MerkleRoot}
\left(
\{H(\mathsf{Encode}(r)):r\in\mathcal{R}_i^{(0)}\}
\right)
\end{equation}
where $\mathcal{R}_i^{(0)}$ denotes its initial revocation set; a
predefined empty-tree root is used when
$\mathcal{R}_i^{(0)}=\emptyset$. The authority then computes
\begin{equation}
C_i^{auth,(0)}
=
H\!\left(
ID_i\parallel Dom_i\parallel AHash_i
\parallel v_i^{(0)}\parallel RevRoot_i^{(0)}
\right)
\end{equation}
Thus, $C_i^{auth,(0)}$ binds the authority identity to its current
attribute namespace, authorization version, and revocation state.

\paragraph*{Step 4: Blockchain Anchoring and Synchronization}

The authenticated authority state
\begin{equation}
State_i^{(0)}
=
\left(
ID_i,v_i^{(0)},AHash_i,
RevRoot_i^{(0)},C_i^{auth,(0)}
\right)
\end{equation}
is anchored on the consortium blockchain and serves as the authoritative
state for subsequent policy-state derivation and freshness verification.
The AIM maintains the current multi-authority version state
\begin{equation}
\mathbf{V}^{cur}
=
\{(ID_i,v_i)\}_{i=1}^{N_{AA}}
\end{equation}
and synchronizes relevant authorization metadata with the corresponding
Fog Search Nodes.

Whenever the attribute namespace or revocation state of $AA_i$ changes,
only that authority advances
\begin{equation}
v_i^{(t+1)}=v_i^{(t)}+1
\end{equation}
and recomputes its authorization-state commitment. Unaffected authorities
retain their existing states, enabling localized authorization updates
without global reinitialization.

\subsubsection*{\textbf{Phase III: User Registration and
Version-Bound Authorization Profile Generation}}

Each Data Owner (DO) or Data User (DU) enrolls with the relevant
Attribute Authorities (AAs) to obtain attribute secret keys. The
Authorization Index Manager (AIM) subsequently constructs a
Version-Bound Authorization Profile (VAP) that binds the user's
authorized attributes and domains to the current states of the issuing
authorities. This profile provides the user-side authorization state for
subsequent policy-state-aware search, synchronization, and verification.

\paragraph*{Step 1: User Registration and Attribute-Key Generation}

User $U$ submits
\begin{equation}
Req_U=(UID,Dom,Role,Cred)
\end{equation}
where $UID$, $Dom$, $Role$, and $Cred$ denote the user identity,
administrative domain, organizational role, and supporting credentials,
respectively. Each relevant $AA_i$ independently verifies $Cred$ and
assigns an authorized attribute subset
\begin{equation}
\mathcal{S}_{U,i}\subseteq\mathcal{A}_i
\end{equation}
The authority then generates the corresponding user-bound attribute key
\begin{equation}
SK_{U,i}
\leftarrow
\mathsf{AA.KeyGen}
(MSK_i,UID,\mathcal{S}_{U,i})
\end{equation}
For the participating authority set $\mathcal{AA}_U$, the user's
attribute-key material is
\begin{equation}
SK_U=
\bigcup_{AA_i\in\mathcal{AA}_U}SK_{U,i}
\end{equation}

\paragraph*{Step 2: Post-Quantum Protected Attribute-Key Delivery}

Using the ML-KEM public key established for $U$ in Phase~I, each
participating $AA_i$ computes
\begin{equation}
(ct_i,ss_i)
\leftarrow
\mathsf{KEM.Encaps}(pk_U^{\mathsf{KEM}})
\end{equation}
and derives
\begin{equation}
K_i=\mathsf{KDF}(ss_i)
\end{equation}
The attribute-key component is protected as
\begin{equation}
EncKey_i=
\mathsf{AE.Enc}(K_i,SK_{U,i})
\end{equation}
and $(ct_i,EncKey_i)$ is delivered to $U$. The user decapsulates $ct_i$
using $sk_U^{\mathsf{KEM}}$, derives $K_i$, and recovers $SK_{U,i}$.
Thus, ML-KEM provides post-quantum protection for attribute-key delivery,
while attribute authorization remains governed by the multi-authority
mechanism.

\paragraph*{Step 3: Version-Bound Authorization Profile Generation}

After successful enrollment, the AIM obtains the authorization metadata
associated with $U$ but does not receive or store the secret key
components $SK_{U,i}$. The user's authorized attribute set is
\begin{equation}
\mathcal{S}_U=
\bigcup_{AA_i\in\mathcal{AA}_U}\mathcal{S}_{U,i}
\end{equation}
From the authoritative states anchored in Phase~II, the AIM obtains the
user-relevant version state
\begin{equation}
\mathbf{V}_U=
\{(ID_i,v_i):AA_i\in\mathcal{AA}_U\}
\end{equation}
and corresponding authority-state commitments
\begin{equation}
\mathcal{C}_U=
\{(ID_i,C_i^{auth}):AA_i\in\mathcal{AA}_U\}
\end{equation}

The AIM computes the user's authorization root as
\begin{equation}
\begin{aligned}
AuthRoot_U =
H\!\Big(&UID
\parallel H(\mathsf{Encode}(\mathcal{S}_U))
\\
&\parallel H(\mathsf{Encode}(\mathbf{V}_U))
\parallel H(\mathsf{Encode}(\mathcal{C}_U))
\Big)
\end{aligned}
\end{equation}
and constructs
\begin{equation}
VAP_U=
\left(
UID,
\mathcal{D}_U,
\mathbf{V}_U,
\mathcal{C}_U,
AuthRoot_U
\right)
\end{equation}
where $\mathcal{D}_U$ denotes the administrative domains authorized for
$U$.

The AIM maintains $VAP_U$ against the latest blockchain-anchored
authority states. When an authority relevant to $U$ updates its
attribute or revocation state, only the corresponding entries in
$\mathbf{V}_U$ and $\mathcal{C}_U$ are refreshed and $AuthRoot_U$ is
recomputed. States of unrelated authorities remain unchanged, enabling
localized authorization synchronization.

\subsubsection*{\textbf{Phase IV: Policy-State-Aware Dynamic Search Index
Construction}}

This phase constructs a verifiable dynamic search index whose entries are
bound to both the access policy and the current states of only the
Attribute Authorities (AAs) governing that policy. This policy-relevant
binding prevents unrelated authority updates from invalidating the index
and enables localized synchronization under policy or revocation changes.
For each encrypted IoMT record already outsourced with identifier $CID_i$,
the Data Owner (DO) constructs the index as follows.

\paragraph*{Step 1: Keyword and Policy-State Extraction}

For record $M_i$, the DO extracts the searchable keyword set
\begin{equation}
W_i=\{w_{i,1},\ldots,w_{i,t}\}
\end{equation}
and associates it with
\begin{equation}
Meta_i=(PID_i,Dom_i,TS_i)
\end{equation}
where $PID_i$ identifies the access policy, $Dom_i$ is the administrative
domain, and $TS_i$ is the timestamp.

Let $\mathcal{AA}(PID_i)$ denote the AAs governing $PID_i$. From the
blockchain-anchored authority states established in Phase~II, the DO
derives the policy-relevant version state
\begin{equation}
\mathbf{V}_{P_i}
=
\{(ID_k,v_k):AA_k\in\mathcal{AA}(PID_i)\}
\label{eq:policy-version-state}
\end{equation}
Thus, $\mathbf{V}_{P_i}$ excludes authority states unrelated to the
record policy.

\paragraph*{Step 2: Policy-State-Bound Keyword Encoding}

The DO first derives a compact policy-state digest
\begin{equation}
PV_i=
H\!\left(
\mathsf{Encode}(\mathbf{V}_{P_i})
\right)
\label{eq:policy-version-digest}
\end{equation}
For each $w_{i,j}\in W_i$, it then computes
\begin{equation}
T_{i,j}
=
H\!\left(
w_{i,j}
\parallel PID_i
\parallel PV_i
\parallel Dom_i
\right)
\label{eq:policy-bound-token}
\end{equation}
Hence, a token is valid for a specific keyword, policy, domain, and
policy-relevant authorization state. An update by an AA outside
$\mathcal{AA}(PID_i)$ does not alter $T_{i,j}$, whereas an affected
policy-state update produces a new token state.

\paragraph*{Step 3: Dynamic Search Index Construction}

Each token is mapped to the outsourced encrypted record through
\begin{equation}
I_{i,j}
=
(T_{i,j},CID_i,PID_i,PV_i)
\label{eq:index-entry}
\end{equation}
and the record-level index is
\begin{equation}
\mathcal{I}_i
=
\{I_{i,j}:w_{i,j}\in W_i\}
\end{equation}
The global Dynamic Search Index (DSI) is maintained as
\begin{equation}
DSI=\bigcup_i\mathcal{I}_i
\end{equation}
Because entries are independently addressable by their policy-state-bound
tokens, affected entries can be inserted, deleted, or replaced without
reconstructing unrelated index partitions.

\paragraph*{Step 4: Authenticated Index Commitment}

To support verifiable retrieval, the DO constructs Merkle leaves
\begin{equation}
L_{i,j}
=
H\!\left(
\mathsf{Encode}(I_{i,j})
\right)
\end{equation}
and computes the record-level index root
\begin{equation}
Root_i
=
\mathsf{MerkleRoot}
\left(
\{L_{i,j}\}_{j=1}^{t}
\right)
\label{eq:index-root}
\end{equation}
The DSI maintains the corresponding authentication paths, allowing a
returned index entry to be verified against $Root_i$ without retrieving
the complete index.

\paragraph*{Step 5: Policy-State-Aware Record Commitment}

The DO derives the authenticated state of the authorities governing
$PID_i$ as
\begin{equation}
AuthState_i
=
H\!\left(
\mathsf{Encode}
\left(
\{(ID_k,C_k^{auth}):
AA_k\in\mathcal{AA}(PID_i)\}
\right)
\right)
\label{eq:policy-auth-state}
\end{equation}
It then computes
\begin{equation}
Commit_i
=
H\!\left(
CID_i
\parallel Root_i
\parallel PID_i
\parallel PV_i
\parallel AuthState_i
\right)
\label{eq:record-commitment}
\end{equation}

The resulting commitment jointly binds the outsourced encrypted-record
identifier, authenticated search index, access policy, and its
policy-relevant multi-authority state. The authenticated metadata
\begin{equation}
AMeta_i=
(CID_i,PID_i,PV_i,Root_i,AuthState_i,Commit_i)
\label{eq:authenticated-metadata}
\end{equation}
is subsequently anchored to the consortium blockchain. Consequently,
updates from unrelated authorities leave $AMeta_i$ unchanged, while
changes affecting $\mathcal{AA}(PID_i)$ require only the corresponding
policy-bound index and commitment state to be incrementally refreshed.


\subsubsection*{\textbf{Phase V: Secure Data Encryption and Outsourcing}}

This phase securely outsources IoMT records while decoupling encrypted
data storage from the policy-state-aware searchable index. Each record is
protected using a fresh symmetric key, while the key is encrypted under
the corresponding multi-authority access policy. The resulting
content-addressed identifier provides an immutable reference to the
current encrypted record state and is subsequently bound to the
searchable index.

\paragraph*{Step 1: Authenticated IoMT Record Encryption}

For each IoMT record $M_i$, the Data Owner (DO) generates a fresh
symmetric data key
\begin{equation}
K_i\overset{\$}{\leftarrow}\{0,1\}^{256}
\end{equation}
Let
\begin{equation}
AAD_i=
\mathsf{Encode}(PID_i,Dom_i,TS_i)
\end{equation}
denote the authenticated metadata associated with the record, where
$PID_i$, $Dom_i$, and $TS_i$ denote the access-policy identifier,
administrative domain, and timestamp, respectively. The DO encrypts
$M_i$ as
\begin{equation}
C_i=
\mathsf{AE.Enc}(K_i,M_i;AAD_i)
\label{eq:data-encryption}
\end{equation}
using a fresh nonce according to the requirements of $\mathsf{AE}$.
Thus, the record and its associated policy metadata are cryptographically
bound by authenticated encryption.

\paragraph*{Step 2: Multi-Authority Policy Protection of Data Key}

Let $P_i$ denote the access policy identified by $PID_i$, and let
$\mathcal{AA}(PID_i)$ be the set of Attribute Authorities governing
$P_i$. The DO protects $K_i$ under $P_i$ as
\begin{equation}
CTK_i=
\mathsf{MAABE.Enc}
\left(
\{PK_k\}_{AA_k\in\mathcal{AA}(PID_i)},
P_i,
K_i
\right)
\label{eq:data-key-protection}
\end{equation}
Only a user possessing valid attribute-key components satisfying $P_i$
can subsequently recover $K_i$. The multi-authority encryption therefore
enforces fine-grained access to the encrypted IoMT record, while the
ML-KEM mechanism of Phase~III separately provides post-quantum protection
for attribute-key delivery.

\paragraph*{Step 3: Encrypted Data Object Construction}

The DO constructs
\begin{equation}
CT_i^{data}
=
(C_i,CTK_i,PID_i,Dom_i,TS_i)
\label{eq:data-object}
\end{equation}
where neither the searchable index nor its integrity commitment is
included in $CT_i^{data}$. This separation allows authorization and
search-index states to evolve independently without re-outsourcing an
unchanged encrypted record.

\paragraph*{Step 4: IPFS Outsourcing and Content Identification}

The DO uploads the encrypted object to IPFS:
\begin{equation}
CID_i=
\mathsf{IPFS.Upload}(CT_i^{data})
\label{eq:ipfs-upload}
\end{equation}
The resulting $CID_i$ provides an immutable content-addressed reference
to the current encrypted record state. The DO maintains
\begin{equation}
DMeta_i=
(CID_i,PID_i,Dom_i,TS_i)
\label{eq:data-metadata}
\end{equation}
for subsequent index construction and blockchain registration. Neither
the plaintext $M_i$ nor the data key $K_i$ is outsourced.

If the encrypted record is subsequently modified, a new encrypted object
$CT_i^{data\prime}$ produces a new identifier $CID_i'$. Only the index
entries and authenticated metadata referring to that record must then be
updated; unchanged encrypted records remain unaffected.

\paragraph*{Step 5: Handoff to Policy-State-Aware Index Construction}

After obtaining $CID_i$, the DO constructs the policy-state-aware
searchable index. The identifier is embedded in each corresponding index
entry, thereby linking searchable metadata to the exact encrypted object
stored in IPFS without embedding the index itself in the content-addressed
ciphertext.

The resulting dependency is
\begin{equation}
M_i
\rightarrow
(C_i,CTK_i)
\rightarrow
CT_i^{data}
\rightarrow
CID_i
\rightarrow
\mathcal{I}_i
\rightarrow
Root_i
\rightarrow
Commit_i
\label{eq:data-index-dependency}
\end{equation}
Consequently, encrypted data storage and searchable-index evolution remain
decoupled: policy or authorization-state changes can update only the
affected index and commitment states, whereas data modifications generate
a new $CID_i$ and require updates only to entries referencing the modified
record.


\subsubsection*{\textbf{Phase VI: Adaptive Authorization-Aware Search Scheduling}}

Upon receiving a search request, the framework first validates the
Data User's (DU's) current authorization state and derives only the
policy-state-aware index shards that the user is permitted to search.
The Adaptive Authorization-Aware Search Scheduler (AASS) then assigns
each required shard to an eligible Fog Search Node (FSN) according to
index locality, authorization-state freshness, verification workload,
and queue delay.

\paragraph*{Step 1: Authorization-Bound Search Request}

Suppose user $U$ searches for
\begin{equation}
Q=\{w_1,\ldots,w_q\}
\end{equation}
Using
\begin{equation}
VAP_U=
(UID,\mathcal{D}_U,\mathbf{V}_U,\mathcal{C}_U,AuthRoot_U)
\end{equation}
the DU submits
\begin{equation}
Req_Q=(Q,UID,AuthRoot_U,\rho)
\label{eq:search-request}
\end{equation}
where $\rho$ is a fresh nonce for replay protection.

\paragraph*{Step 2: Authorization Verification and Token Derivation}

The Authorization Index Manager (AIM) validates $VAP_U$ against the
latest blockchain-anchored authority states. For each policy $P_\ell$
satisfying
\begin{equation}
\mathcal{S}_U\models P_\ell,
\label{eq:policy-satisfaction}
\end{equation}
the AIM verifies that all authorities governing $P_\ell$ are represented
by current entries in $\mathbf{V}_U$ and $\mathcal{C}_U$. It then derives
\begin{equation}
\mathbf{V}_{P_\ell}
=
\{(ID_k,v_k):AA_k\in\mathcal{AA}(PID_\ell)\}
\end{equation}
and the corresponding policy-state digest
\begin{equation}
PV_\ell=
H\!\left(
\mathsf{Encode}(\mathbf{V}_{P_\ell})
\right)
\end{equation}

For each $w_r\in Q$, the AIM computes
\begin{equation}
T_{r,\ell}^{Q}
=
H\!\left(
w_r\parallel PID_\ell\parallel PV_\ell\parallel Dom_\ell
\right)
\label{eq:query-token}
\end{equation}
The resulting token set
\begin{equation}
T_Q=
\{T_{r,\ell}^{Q}\}_{w_r\in Q,\;P_\ell\in\mathcal{P}_U}
\end{equation}
is mapped to the required authorized shard set $\mathcal{S}_Q$, where
$\mathcal{P}_U$ denotes the policies satisfied by $U$. Invalid, stale,
or unauthorized requests are rejected before index traversal.

\paragraph*{Step 3: Adaptive Authorization-Aware Search Scheduling}

For each eligible $FSN_j$, AASS estimates the search cost
\begin{equation}
SC_j=
\lambda_1C_j^{index}
+\lambda_2C_j^{verify}
+\lambda_3C_j^{sync}
+\lambda_4C_j^{queue}
\label{eq:search-cost}
\end{equation}
where
\begin{equation}
C_j^{index}
=|\widehat{Cand}_Q^{(j)}|
\end{equation}
\begin{equation}
C_j^{verify}
&=|\widehat{R}_Q^{(j)}|\log N_j
\end{equation}
\begin{equation}
C_j^{sync}
=
\left|
\{(ID_k,v_k)\in\mathbf{V}_{Q}:v_{j,k}<v_k\}
\right|
\end{equation}
\begin{equation}
C_j^{queue}
=T_j^{queue}
\end{equation}
Here, $\widehat{Cand}_Q^{(j)}$ and $\widehat{R}_Q^{(j)}$ denote the
estimated candidate and result sizes, $N_j$ is the local index size,
$\mathbf{V}_{Q}$ is the query-relevant authority state, and
$T_j^{queue}$ is the current queue delay. In particular,
$C_j^{sync}$ estimates the synchronization overhead from the number of
query-relevant authority states at $FSN_j$ that lag behind their current
versions.

As summarized in Algorithm~\ref{alg:aass}, each required shard is
assigned only to an FSN maintaining that shard, while replicated shards
are routed to the eligible node with the minimum predicted search cost.

\begin{algorithm}[t]
\caption{Adaptive Authorization-Aware Search Scheduling (AASS)}
\label{alg:aass}
\begin{algorithmic}[1]
\Require Token set $T_Q$, required shard set $\mathcal{S}_Q$,
candidate FSNs $\mathcal{F}$
\Ensure Shard-to-FSN assignment $\mathcal{M}_Q$

\State $\mathcal{M}_Q\gets\emptyset$

\ForAll{$S\in\mathcal{S}_Q$}
    \State $FSN^*\gets\bot$, $SC_{\min}\gets\infty$

    \ForAll{$FSN_j\in\mathcal{F}$}
        \If{$S\notin\mathcal{S}_j$}
            \State \textbf{continue}
        \EndIf

        \State $SC_j\gets
        EstimateSearchCost(FSN_j,T_Q,S)$

        \If{$SC_j<SC_{\min}$}
            \State $SC_{\min}\gets SC_j$
            \State $FSN^*\gets FSN_j$
        \EndIf
    \EndFor

    \State $\mathcal{M}_Q[S]\gets FSN^*$
\EndFor

\State Dispatch relevant tokens and shards according to $\mathcal{M}_Q$
\State \Return $\mathcal{M}_Q$

\end{algorithmic}
\end{algorithm}

\paragraph*{Step 4: Distributed Encrypted Search}

Each selected FSN searches only its assigned authorized shard using
$T_Q$. The matching entries are
\begin{equation}
R=
\left\{
I_{i,j}=
(T_{i,j},CID_i,PID_i,PV_i)
:
T_{i,j}\in T_Q
\right\}
\label{eq:search-results}
\end{equation}
Because~\eqref{eq:query-token} follows the same policy-state encoding
used during index construction, a matching entry is bound to the
validated policy and its current policy-relevant authority state.

\paragraph*{Step 5: Verifiable Search Response Generation}

For each returned entry $I_{i,j}$, the responsible FSN obtains its
Merkle authentication path $\pi_{i,j}$ and generates
\begin{equation}
\Pi_{i,j}
=
(I_{i,j},Root_i,AuthState_i,Commit_i,\pi_{i,j})
\label{eq:verification-bundle}
\end{equation}
where $\pi_{i,j}$ authenticates the leaf
\begin{equation}
L_{i,j}
=
H\!\left(
\mathsf{Encode}(I_{i,j})
\right)
\end{equation}
against $Root_i$.

The aggregated response
\begin{equation}
Resp=
\{\Pi_{i,j}\}_{I_{i,j}\in R}
\label{eq:search-response}
\end{equation}
is returned to the DU for integrity and authorization-state verification
before the corresponding encrypted objects are retrieved from IPFS.

\subsubsection*{\textbf{Phase VII: Dependency-Aware Dynamic Evolution
and Incremental Authorization Synchronization}}

Dynamic IoMT environments require efficient handling of data changes,
policy updates, and revocations without rebuilding the global searchable
index or synchronizing all Fog Search Nodes (FSNs). We therefore introduce
Dependency-Aware Incremental Authorization Synchronization (DIAS), which
identifies the dependency closure of each update and evolves only the
affected policy, index, user-profile, and FSN states.

\paragraph*{Step 1: Update Classification and Dependency Discovery}

An update is represented as
\begin{equation}
U=(Op,\mathcal{X},TS_U)
\end{equation}
where
\begin{equation}
Op\in\{\mathsf{Insert},\mathsf{Modify},\mathsf{Delete},
\mathsf{PolicyUpdate},\mathsf{Revoke}\}
\end{equation}
For a state change at $AA_k$, the AIM identifies
\begin{equation}
\mathcal{P}_k^{aff}
=\{P_\ell:AA_k\in\mathcal{AA}(PID_\ell)\}
\end{equation}
\begin{equation}
\mathcal{R}_k^{aff}
=\{i:P_i\in\mathcal{P}_k^{aff}\}
\end{equation}
\begin{equation}
\mathcal{S}_k^{aff}
=\{S:I_{i,j}\in S,\ i\in\mathcal{R}_k^{aff}\}
\label{eq:dependency-closure}
\end{equation}
Thus, an authorization change propagates only through policies, records,
and index shards that depend on the affected authority.

\paragraph*{Step 2: Localized Policy-State and Index Evolution}

When $AA_k$ changes its authorization state, it advances
$v_k'=v_k+1$ and publishes $C_k^{auth\prime}$ as in Phase~II.
For each affected policy $P_i$, the AIM derives
\begin{equation}
\mathbf{V}_{P_i}'
=
\{(ID_h,v_h'):AA_h\in\mathcal{AA}(PID_i)\}
\end{equation}
\begin{equation}
PV_i'=H(\mathsf{Encode}(\mathbf{V}_{P_i}'))
\label{eq:updated-policy-state}
\end{equation}
The affected tokens and index entries are updated as
\begin{equation}
T_{i,j}'
=
H(w_{i,j}\parallel PID_i\parallel PV_i'\parallel Dom_i)
\end{equation}
\begin{equation}
I_{i,j}'
=
(T_{i,j}',CID_i,PID_i,PV_i')
\label{eq:updated-index}
\end{equation}
Accordingly,
\begin{equation}
\Delta\mathcal{I}_i
=
\{(I_{i,j},I_{i,j}')\}_{j\in J_i^{aff}}
\end{equation}
contains only modified entries. In particular,
\begin{equation}
AA_k\notin\mathcal{AA}(PID_i)
\Longrightarrow
PV_i'=PV_i
\label{eq:unaffected-policy}
\end{equation}
so unrelated authority updates do not trigger index evolution.

\paragraph*{Step 3: Incremental Authenticated-State Evolution}

Only affected Merkle leaves are replaced:
\begin{equation}
Root_i'
=
\mathsf{MerkleUpdate}(Root_i,\Delta\mathcal{L}_i)
\end{equation}
where
$L_{i,j}'=H(\mathsf{Encode}(I_{i,j}'))$.
The updated policy-authority state and commitment are
\begin{align}
AuthState_i'
&=
H\!\left(
\mathsf{Encode}
\{(ID_h,C_h^{auth\prime}):
AA_h\in\mathcal{AA}(PID_i)\}
\right)\\
Commit_i'
&=
H(CID_i\parallel Root_i'\parallel PID_i
\parallel PV_i'\parallel AuthState_i')
\label{eq:updated-commitment}
\end{align}
For a data modification, Phase~V instead produces a new $CID_i'$, which
replaces $CID_i$ in the affected index entries and commitment.

\paragraph*{Step 4: Dependency-Aware Incremental Synchronization}

The AIM constructs the compact update
\begin{equation}
DIAS_i=
(\Delta\mathbf{V}_{P_i},
\Delta\mathcal{I}_i,
Root_i',PV_i',AuthState_i',Commit_i')
\label{eq:dias-message}
\end{equation}
and sends it only to
\begin{equation}
\mathcal{F}_i^{aff}
=
\{FSN_j:\mathcal{S}_j\cap\mathcal{S}_i^{aff}\neq\emptyset\}
\end{equation}
Each receiving FSN updates its affected shards and corresponding
components of $\mathbf{V}_j$. Likewise, only users satisfying
$AA_k\in\mathcal{AA}_U$ refresh the affected entries of
$\mathbf{V}_U$ and $\mathcal{C}_U$ and recompute $AuthRoot_U$.

\paragraph*{Step 5: Version-Gated Activation and Anchoring}

An updated shard becomes active only after its required DIAS state has
been applied. Hence, an FSN satisfying
\begin{equation}
\exists(ID_k,v_k)\in\mathbf{V}_Q:\;v_{j,k}<v_k
\end{equation}
remains stale for that query state and is penalized through
$C_j^{sync}$ in Phase~VI.

Finally, the compact update state
\begin{equation}
BC_i'
=
(CID_i',PID_i,PV_i',Root_i',AuthState_i',Commit_i',TS_i')
\label{eq:update-anchor}
\end{equation}
is anchored to the consortium blockchain, where $CID_i'=CID_i$ when
the encrypted record is unchanged. DIAS therefore confines update and
synchronization overhead to the dependency chain
\begin{equation}
AA_k
\rightarrow
\mathcal{P}_k^{aff}
\rightarrow
\mathcal{R}_k^{aff}
\rightarrow
\mathcal{S}_k^{aff}
\rightarrow
\mathcal{F}_k^{aff}
\label{eq:dias-dependency-chain}
\end{equation}
rather than globally rebuilding or redistributing authorization and
search states.


\subsubsection*{\textbf{Phase VIII: Verifiable Retrieval and Secure Decryption}}

After receiving the search response, the Data User (DU) verifies each
returned index entry against its authenticated index and the latest
blockchain-anchored policy state before retrieving and decrypting the
corresponding IoMT record.

\paragraph*{Step 1: Search-Result Integrity Verification}

For each returned
\begin{equation}
\Pi_{i,j}
=
(I_{i,j},Root_i,AuthState_i,Commit_i,\pi_{i,j})
\end{equation}
the DU computes
\begin{equation}
L_{i,j}
=
H\!\left(\mathsf{Encode}(I_{i,j})\right)
\end{equation}
and verifies
\begin{equation}
\mathsf{VerifyMerkle}(Root_i,L_{i,j},\pi_{i,j})=1
\label{eq:result-verification}
\end{equation}
A failed proof causes the corresponding result to be rejected.

\paragraph*{Step 2: Policy-State and Blockchain Verification}

The DU obtains the latest authenticated metadata
\begin{equation}
AMeta_i=
(CID_i,PID_i,PV_i,Root_i,AuthState_i,Commit_i)
\end{equation}
from the blockchain and verifies that the returned values match the
anchored state. It then recomputes
\begin{equation}
Commit_i^*
=
H\!\left(
CID_i\parallel Root_i\parallel PID_i
\parallel PV_i\parallel AuthState_i
\right)
\end{equation}
and accepts the result only if
\begin{equation}
Commit_i^*=Commit_i
\end{equation}

The DU additionally verifies that the authority components governing
$PID_i$ are current in $\mathbf{V}_U$ and that
\begin{equation}
\mathcal{S}_U\models P_i
\label{eq:retrieval-authorization}
\end{equation}
Thus, an authenticated but stale or unauthorized result is rejected
before data retrieval.

\paragraph*{Step 3: Encrypted Record Retrieval}

After successful verification, the DU retrieves
\begin{equation}
CT_i^{data}
=
\mathsf{IPFS.Get}(CID_i)
\end{equation}
and parses
\begin{equation}
CT_i^{data}
=
(C_i,CTK_i,PID_i,Dom_i,TS_i)
\end{equation}
Because $CID_i$ is included in the blockchain-anchored commitment, the
retrieved encrypted object is bound to the verified search result.

\paragraph*{Step 4: Multi-Authority Key Recovery and Decryption}

If the user's attributes satisfy $P_i$, the DU recovers
\begin{equation}
K_i=
\mathsf{MAABE.Dec}(SK_U,CTK_i)
\end{equation}
and reconstructs
\begin{equation}
AAD_i=
\mathsf{Encode}(PID_i,Dom_i,TS_i)
\end{equation}
The IoMT record is recovered by
\begin{equation}
M_i=
\mathsf{AE.Dec}(K_i,C_i;AAD_i)
\label{eq:final-decryption}
\end{equation}
The record is accepted only if both the multi-authority policy check and
authenticated decryption succeed.

\section{Security Analysis}
\label{sec:security}

This section analyzes MA-LB-PQ-VDSE with respect to outsourced-data
confidentiality, authorization-state consistency, policy-state-aware
search, dynamic-update integrity, verifiable retrieval, and
post-quantum protected attribute-key delivery.

\subsection{Security Model and Assumptions}

We consider PPT adversaries controlling IPFS or one or more Fog Search
Nodes (FSNs). Such adversaries may observe stored ciphertexts and index
entries, replay stale states, omit or modify returned entries, and attempt
to forge verification evidence. A revoked user may retain previously
issued credentials and attempt to reuse stale authorization or search
state.

The Attribute Authorities (AAs) and Data Owners (DOs) are trusted to
execute the protocol correctly. The Authorization Index Manager (AIM) is
trusted for authorization processing: it maintains user authorization
metadata, validates VAPs, derives policy-state-aware query tokens, and
coordinates DIAS updates, but it does not obtain users' attribute secret
keys or plaintext IoMT records. The consortium blockchain is assumed to
provide authenticated consensus and immutable finalized state.

The security analysis relies on: (i) collision and preimage resistance
of $H$; (ii) IND-CPA security of the underlying multi-authority ABE
scheme; (iii) confidentiality and ciphertext integrity of
$\mathsf{AE}$; (iv) IND-CCA security of ML-KEM against QPT adversaries;
(v) Merkle-tree integrity under collision resistance of $H$; and
(vi) integrity and immutability of finalized blockchain metadata.
Accordingly, the post-quantum guarantee of the current construction is
limited to ML-KEM-protected attribute-key delivery; the pairing-based
multi-authority authorization mechanism is not claimed to be
post-quantum secure.

\subsection{Leakage Model}

The untrusted storage and search infrastructure may observe structural
search leakage. We model
\begin{equation}
\mathcal{L}
=
(\mathcal{L}_{setup},
 \mathcal{L}_{query},
 \mathcal{L}_{update},
 \mathcal{L}_{access})
\end{equation}
where
\begin{equation}
\mathcal{L}_{setup}
=(n,m,\mathcal{D})
\end{equation}
\begin{equation}
\mathcal{L}_{query}
=(|T_Q|,|\mathcal{S}_Q|,|R|)
\end{equation}
\begin{equation}
\mathcal{L}_{update}
=(Op,|\Delta\mathcal{I}|,|\mathcal{F}^{aff}|)
\end{equation}
\begin{equation}
\mathcal{L}_{access}
=\mathsf{Acc}(Q)
\end{equation}
Here, $n$ and $m$ denote the numbers of encrypted records and index
entries, $|T_Q|$ is the number of query tokens, $|\mathcal{S}_Q|$ the
number of accessed shards, $|R|$ the result size, and
$|\mathcal{F}^{aff}|$ the number of FSNs affected by a DIAS update.
$\mathsf{Acc}(Q)$ captures the identifiers of records returned by the
query.

The AIM necessarily learns authorization metadata and the query
information required to evaluate $\mathcal{S}_U\models P_i$ and derive
$T_Q$; these values are therefore outside the privacy boundary with
respect to AIM. Plaintext IoMT records and user attribute secret keys
remain hidden from AIM, FSNs, IPFS, and external observers.


\subsection{Security Theorems}

\begin{theorem}[Confidentiality of Outsourced IoMT Data]
Assume that $\mathsf{AE}$ provides authenticated confidentiality and the
underlying multi-authority ABE scheme is IND-CPA secure. Then an
adversary that does not possess attribute keys satisfying $P_i$ has only
negligible advantage in distinguishing the plaintext $M_i$ protected by
$CT_i^{data}$.
\end{theorem}

\begin{proof}
For each record, the DO independently samples a fresh symmetric key
$K_i$ and computes
\begin{equation}
C_i=\mathsf{AE.Enc}(K_i,M_i;AAD_i)
\end{equation}
while $K_i$ is protected under policy $P_i$ as
\begin{equation}
CTK_i=
\mathsf{MAABE.Enc}
(\{PK_k\}_{AA_k\in\mathcal{AA}(PID_i)},P_i,K_i)
\end{equation}
Consider an adversary that observes
$CT_i^{data}=(C_i,CTK_i,PID_i,Dom_i,TS_i)$ but does not possess
attribute-key components satisfying $P_i$. By IND-CPA security of
$\mathsf{MAABE}$, the adversary cannot distinguish or recover the protected $K_i$ except with negligible probability.

Conditioned on $K_i$ remaining hidden, $C_i$ is an authenticated
encryption of $M_i$ under an independently random key. Therefore, any non-negligible advantage in distinguishing two equal-length candidate
plaintexts from $C_i$ can be converted into an adversary against the
confidentiality of $\mathsf{AE}$. Moreover, modification of $C_i$ or its
associated data is detected by authenticated decryption.

Hence, by a standard hybrid argument,
\begin{equation}
Adv_{\mathsf{data}}
\le
Adv_{\mathsf{MAABE}}
+
Adv_{\mathsf{AE}}
=
\mathsf{negl}(\lambda)
\end{equation}
Thus, an unauthorized adversary learns no plaintext information beyond
the explicitly permitted leakage.
\end{proof}


\begin{theorem}[Policy-State Token Consistency]
For an authorized query keyword $w_r$ and policy $P_\ell$, the query
token matches an index token if both are generated for the same keyword,
policy, domain, and policy-relevant authority state.
\end{theorem}

\begin{proof}
During index construction, Phase~IV computes
\begin{equation}
T_{i,j}
=
H(w_{i,j}\parallel PID_i\parallel PV_i\parallel Dom_i)
\end{equation}
whereas Phase~VI derives
\begin{equation}
T_{r,\ell}^{Q}
=
H(w_r\parallel PID_\ell\parallel PV_\ell\parallel Dom_\ell)
\end{equation}
If
\begin{equation}
w_r=w_{i,j}
\end{equation}
\begin{equation}
PID_\ell=PID_i
\end{equation}
\begin{equation}
PV_\ell=PV_i
\end{equation}
\begin{equation}
Dom_\ell=Dom_i
\end{equation}
both computations receive exactly the same canonically ordered input.
Since $H$ is deterministic,
\begin{equation}
T_{r,\ell}^{Q}=T_{i,j}
\end{equation}
and the corresponding index entry can therefore be matched.

Conversely, suppose an authority governing $P_\ell$ advances its state.
Then the encoded policy-relevant version vector changes, so
\begin{equation}
\mathbf{V}_{P_\ell}'\neq\mathbf{V}_{P_\ell}
\end{equation}
Consequently,
\begin{equation}
PV_\ell'
=
H(\mathsf{Encode}(\mathbf{V}_{P_\ell}'))
\neq
PV_\ell
\end{equation}
except with negligible probability of a hash collision. The resulting
token is therefore bound to the updated policy state. Hence, successful
token matching occurs only for a consistent keyword, policy, domain, and
policy-relevant authorization state, except with negligible collision
probability.
\end{proof}


\begin{theorem}[Correctness of AASS]
For every required shard $S\in\mathcal{S}_Q$, Algorithm~\ref{alg:aass}
selects an eligible FSN with minimum predicted search cost among all FSNs
that maintain $S$, provided at least one such FSN exists.
\end{theorem}

\begin{proof}
For a required shard $S$, Algorithm~\ref{alg:aass} examines the candidate
set
\begin{equation}
\mathcal{F}_S=
\{FSN_j\in\mathcal{F}:S\in\mathcal{S}_j\}
\end{equation}
Nodes outside $\mathcal{F}_S$ are skipped and therefore cannot be selected
for a shard they do not maintain. By assumption,
$\mathcal{F}_S\neq\emptyset$.

For every $FSN_j\in\mathcal{F}_S$, AASS evaluates
\begin{equation}
SC_j=
\lambda_1C_j^{index}
+\lambda_2C_j^{verify}
+\lambda_3C_j^{sync}
+\lambda_4C_j^{queue}
\end{equation}
using the same cost model. The algorithm initializes
$SC_{\min}=\infty$ and replaces the current choice whenever a lower value
is observed. After all eligible nodes are examined, the stored node
therefore satisfies
\begin{equation}
FSN_S^*
=
\arg\min_{FSN_j:S\in\mathcal{S}_j}SC_j
\end{equation}
Thus, the selected node is both feasible and minimum-cost under the
defined prediction model.

Applying the same procedure independently to every
$S\in\mathcal{S}_Q$ produces the shard-to-FSN map $\mathcal{M}_Q$ in
which every required shard is assigned to an eligible minimum-cost FSN.
Hence AASS is correct with respect to its stated scheduling objective.
\end{proof}


\begin{theorem}[Policy-State Non-Interference]
Let $AA_k$ update its authorization state. For any record $i$ satisfying
$AA_k\notin\mathcal{AA}(PID_i)$, the update does not change its
policy-state digest or searchable-index entries, except with negligible
probability of a hash collision.
\end{theorem}

\begin{proof}
For record $i$, the policy-relevant authority state is defined only over
the authorities governing $PID_i$:
\begin{equation}
\mathbf{V}_{P_i}
=
\{(ID_h,v_h):AA_h\in\mathcal{AA}(PID_i)\}
\end{equation}
If
\begin{equation}
AA_k\notin\mathcal{AA}(PID_i)
\end{equation}
then advancing $v_k$ does not modify any component of
$\mathbf{V}_{P_i}$. Therefore,
\begin{equation}
\mathbf{V}_{P_i}'=\mathbf{V}_{P_i}
\end{equation}
and, by deterministic canonical encoding,
\begin{equation}
PV_i'
=
H(\mathsf{Encode}(\mathbf{V}_{P_i}'))
=
H(\mathsf{Encode}(\mathbf{V}_{P_i}))
=
PV_i
\end{equation}

Since $PID_i$, $Dom_i$, and the indexed keyword $w_{i,j}$ are also
unchanged, the token remains
\begin{equation}
T_{i,j}'
=
H(w_{i,j}\parallel PID_i\parallel PV_i'\parallel Dom_i)
=
T_{i,j}
\end{equation}
Consequently,
\begin{equation}
I_{i,j}'=(T_{i,j}',CID_i,PID_i,PV_i')
=
I_{i,j}
\end{equation}
Hence no index delta is generated for record $i$. Thus, an authority
update propagates only to records whose policies actually depend on that
authority, establishing policy-state non-interference.
\end{proof}


\begin{theorem}[DIAS State Consistency and Stale-State Exclusion]
Assume collision resistance of $H$ and immutable finalized blockchain
state. An FSN whose query-relevant authority state is older than the
current DIAS state cannot be treated as current for that policy, and a
tampered affected index state cannot produce the anchored commitment
except with negligible probability.
\end{theorem}

\begin{proof}
For each affected policy,DIAS derives the updated state comprising $PV_i'$, $\Delta\mathcal{I}_i$, $Root_i'$, $AuthState_i'$, and $Commit_i'$.

The resulting commitment is
\begin{equation}
Commit_i'
=
H(CID_i'\parallel Root_i'\parallel PID_i
\parallel PV_i'\parallel AuthState_i')
\end{equation}
Thus, $Commit_i'$ jointly binds the current encrypted-object reference,
authenticated index root, policy identifier, policy-version digest, and
policy-relevant authority commitments.

Suppose an adversary modifies an affected index entry, substitutes a
different $CID_i'$, changes $PV_i'$ or $AuthState_i'$, or replaces
$Root_i'$. The recomputed commitment will differ from the
blockchain-anchored $Commit_i'$ unless the adversary finds two distinct
encodings producing the same hash. By collision resistance of $H$, this
occurs only with negligible probability. Replacing the anchored
commitment itself would contradict the assumed integrity of finalized
blockchain state.

Freshness is enforced separately by version-gated activation. Phase~VII
does not treat an updated shard as current until the required DIAS state
has been applied. In particular, if
\begin{equation}
\exists(ID_k,v_k)\in\mathbf{V}_Q:\;v_{j,k}<v_k
\end{equation}
then $FSN_j$ is explicitly identified as stale for that query-relevant
authority state and is accounted for through $C_j^{sync}$ before search
execution.

Therefore, neither a stale state nor a tampered affected state can be
accepted as the current DIAS state except with negligible probability.
\end{proof}


\begin{theorem}[Integrity and Freshness of Verifiable Retrieval]
Assume collision resistance of $H$, Merkle-tree integrity, and immutable
finalized blockchain metadata. A malicious FSN cannot cause a forged,
modified, or stale returned index entry to be accepted by a legitimate DU
except with negligible probability.
\end{theorem}

\begin{proof}
For each returned entry $I_{i,j}$, the DU first reconstructs the exact
Merkle leaf
\begin{equation}
L_{i,j}=H(\mathsf{Encode}(I_{i,j}))
\end{equation}
and accepts its membership only if
\begin{equation}
\mathsf{VerifyMerkle}(Root_i,L_{i,j},\pi_{i,j})=1
\end{equation}
Hence, replacing or modifying any component of $I_{i,j}$ while retaining
a valid authentication path would require forging a Merkle proof or
finding a collision in $H$.

The DU then obtains the latest blockchain-anchored metadata
\begin{equation}
AMeta_i=
(CID_i,PID_i,PV_i,Root_i,AuthState_i,Commit_i)
\end{equation}
and recomputes
\begin{equation}
Commit_i^*
=
H(CID_i\parallel Root_i\parallel PID_i
\parallel PV_i\parallel AuthState_i)
\end{equation}
Acceptance requires $Commit_i^*=Commit_i$ and agreement with the latest
anchored metadata. Therefore, substitution of the ciphertext reference,
Merkle root, policy identifier, policy-state digest, or authority-state
commitment without detection would imply either a collision in $H$ or a
violation of blockchain immutability.

Finally, the DU verifies that the authority components governing
$PID_i$ are current in $\mathbf{V}_U$ and that
\begin{equation}
\mathcal{S}_U\models P_i
\end{equation}
Thus, a previously valid but stale index state cannot be accepted merely
because its old Merkle proof remains internally valid.

Accordingly, forging, modifying, or replaying an accepted returned entry
requires breaking Merkle integrity, hash collision resistance, or the
integrity of the latest blockchain state. Under the stated assumptions,
the success probability is negligible.
\end{proof}


\begin{theorem}[Post-Quantum Protection of Attribute-Key Delivery]
Assume ML-KEM is IND-CCA secure against QPT adversaries, $\mathsf{KDF}$
is secure, and $\mathsf{AE}$ provides authenticated confidentiality.
Then a QPT network adversary cannot recover an attribute-key component
$SK_{U,i}$ from $(ct_i,EncKey_i)$ except with negligible probability.
\end{theorem}

\begin{proof}
During Phase~III, authority $AA_i$ encapsulates a fresh shared secret as
\begin{equation}
(ct_i,ss_i)
\leftarrow
\mathsf{KEM.Encaps}(pk_U^{\mathsf{KEM}})
\end{equation}
and derives
\begin{equation}
K_i=\mathsf{KDF}(ss_i)
\end{equation}
The user's attribute-key component is then protected as
\begin{equation}
EncKey_i=
\mathsf{AE.Enc}(K_i,SK_{U,i})
\end{equation}
A network adversary observes $(ct_i,EncKey_i)$ but does not possess
$sk_U^{\mathsf{KEM}}$.

Consider first replacing $ss_i$ obtained from ML-KEM with an independent
uniform secret. A QPT adversary that distinguishes this replacement with
non-negligible advantage yields an adversary against the IND-CCA security
of ML-KEM. Given an indistinguishable shared secret, security of
$\mathsf{KDF}$ implies that the derived $K_i$ is computationally
indistinguishable from an independent encryption key.

Conditioned on $K_i$ being hidden, recovery or distinguishing of
$SK_{U,i}$ from $EncKey_i$ contradicts the authenticated confidentiality
of $\mathsf{AE}$. Consequently,
\begin{equation}
Adv_{\mathsf{PQKey}}
\le
Adv_{\mathsf{ML\mbox{-}KEM}}^{\mathsf{QPT}}
+
Adv_{\mathsf{KDF}}
+
Adv_{\mathsf{AE}}
=
\mathsf{negl}(\lambda)
\end{equation}
Therefore, attribute-key delivery remains protected against the specified
QPT network adversary under the stated assumptions.
\end{proof}




\section{Evaluation}
\label{sec:evaluation}

This section evaluates MA-LB-PQ-VDSE through analytical computation-cost
comparison and experimental performance analysis. Four representative
schemes are used as baselines: Peony++~\cite{ref30} for verifiable
multilevel DSSE with fine-grained access control and public verification;
Scheme~\cite{ref35} for verifiable dynamic searchable encryption;
Scheme~\cite{ref41} for multi-authority attribute-based searchable
encryption; and Scheme~\cite{ref54} for post-quantum attribute-based
searchable encryption with hybrid indexing and verifiable retrieval.
Together, they cover complementary capabilities in dynamic and verifiable
search, multi-authority authorization, forward/backward privacy,
post-quantum protection, and hybrid-index verification.
\subsection{Computation Cost Analysis}

Table~\ref{tab:cost} compares the dominant computation costs of the
proposed framework with four representative baselines. We consider
search-token generation, encrypted search, dynamic index update,
result verification, and authorization synchronization. N/A indicates
that the corresponding operation is not explicitly supported.

Scheme~\cite{ref30} uses constrained-PRF token generation, with search
and public verification scaling as $O(n_w^l)$ and update cost $O(k)$.
The remaining baseline costs follow their dominant cryptographic and
index-processing operations.

\begin{table}[t]
\centering
\caption{Notation used in the computation-cost analysis.}
\label{tab:cost-notation}
\renewcommand{\arraystretch}{1.05}
\setlength{\tabcolsep}{4pt}
\begin{tabular}{ll}
\hline
\textbf{Notation} & \textbf{Description} \\
\hline
$q$ & Number of queried keywords \\
$|\mathcal{T}_Q|$ & Number of generated query tokens \\
$r$ & Number of returned index entries \\
$k$ & Number of updated index entries/keywords \\
$t$ & Number of keywords in a record-level index \\
$u,e$ & User attributes and satisfying LSSS rows \\
$x$ & Number of candidate/returned documents \\
$n_w^l$ & Matching files for keyword $w$ at level $l$ \\
$a$ & Number of changed policy-relevant authority components \\
$N$ & Total indexed records/ciphertexts \\
\hline
$T_H$ & Hash computation \\
$T_P,T_E$ & Pairing and exponentiation/encryption \\
$T_{\rm PRF},T_{\rm CPRF}$ & PRF and constrained-PRF operations \\
$T_L$ & Encrypted-index lookup/set operation \\
$T_{\rm Punc}$ & Puncturable-PRF operation \\
$T_{\rm MAC},T_{\rm Ver}$ & MAC and signature verification \\
$T_{\rm BC}$ & Blockchain state read \\
$T_{\rm Mul}$ & Multiplication operation  \\
$T_{\rm Sig}$ & Signature generation operation \\
$T_{\rm BF}$  & Bitmap-filter check operation \\
\hline
\end{tabular}
\end{table}

\begin{table*}[t]
\centering
\caption{Dominant computation-cost comparison for searchable retrieval.}
\label{tab:cost}
\renewcommand{\arraystretch}{1.15}
\setlength{\tabcolsep}{3pt}
\begin{tabular}{lccccc}
\hline
\textbf{Scheme} &
\textbf{Token} &
\textbf{Search} &
\textbf{Dynamic} &
\textbf{Verification} &
\textbf{Authorization} \\
&
\textbf{Generation} &
&
\textbf{Update} &
&
\textbf{Synchronization} \\
\hline

Scheme~\cite{ref30}
&
$O(q)T_{\rm CPRF}$
&
$O(n_w^l)$
&
$O(k)$
&
$O(n_w^l)$
&
N/A
\\

Scheme~\cite{ref35}
&
$O(1)T_{\rm PRF}$
&
$O(xq)$
&
\makecell{
$O(t)(T_H+T_{\rm PRF})$\\
$+T_{\rm Punc}$
}
&
$O(x)T_{\rm MAC}$
&
N/A
\\

Scheme~\cite{ref41}
&
\makecell{
$O(u+q)(T_H+T_{\rm Mul})$\\
$+T_E$
}
&
\makecell{
$(2e+1)T_P$\\
$+O(e)(T_E+T_{\rm Mul})$
}
&
N/A
&
N/A
&
N/A
\\

Scheme~\cite{ref54}
&
$O(|\mathcal{T}_Q|)T_{\rm PRF}+T_{\rm Sig}$
&
\makecell{
$O(|\mathcal{T}_Q|)T_L$\\
$+O(x)T_{\rm BF}$
}
&
N/A
&
\makecell{
$O(x\log N)T_H$\\
$+O(x)(T_{\rm Ver}+T_{\rm MAC})$
}
&
N/A
\\

\textbf{Proposed}
&
$O(|\mathcal{T}_Q|)T_H$
&
$O(|\mathcal{T}_Q|)T_L$
&
$O(k\log t)T_H$
&
\makecell{
$O(r\log t)T_H$\\
$+O(r)T_{\rm BC}$
}
&
$O(a+k\log t)T_H$
\\
\hline
\end{tabular}
\end{table*}

For the proposed framework, token generation scales with
$|\mathcal{T}_Q|$, since Phase~VI generates policy-state-aware tokens
for the queried keywords under the authorized policies. Search then
performs direct token lookups only over the shards selected by AASS,
giving $O(|\mathcal{T}_Q|)T_L$ rather than scanning the global index.

Dynamic-update cost is localized to the $k$ affected entries. Updating
their policy-state-bound tokens and record-level Merkle authentication
state requires at most $O(k\log t)T_H$, where $t$ is the number of
indexed keywords for the affected record. Likewise, verification of
$r$ returned entries requires $O(r\log t)T_H$ for their Merkle paths
and $O(r)T_{\rm BC}$ blockchain-state reads.

DIAS synchronization depends only on the changed policy-relevant
authority components and affected index entries, yielding
$O(a+k\log t)T_H$. Unrelated policies, records, and FSNs incur no update
cost. Thus, the proposed framework shifts online processing toward
policy-scoped hash and index operations while avoiding global index
reconstruction and system-wide authorization synchronization.
\subsection{Performance Evaluation}

\subsubsection{Experimental Setup}

The proposed framework was implemented in Python~3.11 and evaluated
on Amazon AWS EC2 c6i.xlarge instances, each equipped with 4~vCPUs
and 8~GB RAM. Four Fog Search Nodes (FSNs) and one cloud server were
deployed within an AWS Virtual Private Cloud (VPC) to emulate a
distributed cross-domain IoMT healthcare environment. Each FSN
maintains an authorization-aware searchable index and performs
encrypted search for its administrative domain, while the cloud
server stores encrypted EHRs and associated blockchain metadata.

The implementation employs ML-KEM-768 for post-quantum key
establishment, AES-256-GCM for symmetric encryption, SHA-256 for
hashing and Merkle-tree construction, and Hyperledger Fabric~v2.5
for maintaining authorization commitments and audit records.
Search indexes combine bitmap-based authorization-aware filtering
with authenticated Merkle trees to support efficient encrypted
retrieval and result verification.

Experiments were conducted using the Synthea synthetic healthcare
dataset. Encrypted searchable indexes containing $10^4$--$10^6$ EHRs were generated from the dataset and uniformly
distributed across four administrative healthcare domains. Unless
otherwise specified, each query contains five keywords. Each
experiment was repeated 10 times, and the reported results represent
the mean with 95\% confidence intervals.

The proposed framework was compared with four representative
searchable encryption schemes: Scheme~\cite{ref30} (Peony++),
Scheme~\cite{ref35}, Scheme~\cite{ref41}, and
Scheme~\cite{ref54} (LV-PQ-ABSE). The evaluation examines trapdoor
generation, encrypted-search efficiency, cross-domain search
scalability, keyword-update performance, result verification,
authorization synchronization, and authorization-aware load
balancing under increasing workloads.


\subsubsection*{\textbf{Experiment 1: Policy-State-Aware Token Generation}}
\label{exp:trapdoor}

This experiment evaluates the online cost of generating policy-state-aware
search tokens. The number of queried keywords is varied as
$q\in\{1,5,10,15,20\}$, while the number of authorized policies is varied
as $|\mathcal{P}_U|\in\{1,2,4,8\}$. The proposed framework is compared
with Guo \textit{et al}.~\cite{ref35}, XB-Muse~\cite{ref36}, Thingom
\textit{et al}.~\cite{ref41}, and Zhuang \textit{et al}.~\cite{ref52}.
ML-KEM key establishment is excluded because it is performed during
attribute-key delivery rather than for each search request.

For every authorized policy $P_\ell\in\mathcal{P}_U$, the AIM derives
the current policy-state digest
\begin{equation}
PV_\ell=H(\mathsf{Encode}(\mathbf{V}_{P_\ell}))
\end{equation}
and generates, for each queried keyword $w_r$,
\begin{equation}
T_{r,\ell}^{Q}
=
H(w_r\parallel PID_\ell\parallel PV_\ell\parallel Dom_\ell)
\end{equation}
Hence, the number of generated tokens is
\begin{equation}
|\mathcal{T}_Q|
=
q|\mathcal{P}_U|
\end{equation}
when all queried keywords are evaluated under all authorized policies,
and the dominant token-generation cost is
\[
O(|\mathcal{T}_Q|)T_H
\]

\begin{figure}[!t]
    \centering
    \includegraphics[width=\columnwidth]{images/fig_exp1_trapdoor.pdf}
    \caption{Policy-state-aware token-generation latency under increasing
    query and authorization scopes.}
    \label{fig:exp1}
\end{figure}

Fig.~\ref{fig:exp1} shows that the proposed scheme exhibits approximately
linear growth with the number of generated policy-aware tokens. Increasing
$q$ enlarges the keyword dimension of the query, while increasing
$|\mathcal{P}_U|$ expands the number of policy states under which those
keywords must be encoded. Nevertheless, token generation remains
hash-dominated and avoids pairing, exponentiation, or lattice operations
in the online search path.

Compared with schemes whose trapdoor construction involves
attribute-dependent group operations or more expensive cryptographic
processing, MA-LB-PQ-VDSE maintains relatively low online preparation
cost. More importantly, the additional computation directly reflects
the authorization scope: only policies satisfying
$\mathcal{S}_U\models P_\ell$ contribute tokens. Thus, token-generation
overhead scales with the effective authorized search scope rather than
the total number of policies or domains in the system.


\subsubsection*{\textbf{Experiment 2: Search Latency}}
\label{exp:search}

This experiment evaluates encrypted-search scalability as the searchable
index size increases from $10^4$ to $10^6$ encrypted records. The
proposed framework is compared with Scheme~\cite{ref30},
Scheme~\cite{ref35}, Scheme~\cite{ref41}, and
Scheme~\cite{ref54}. The query size, number of participating domains,
and authorization scope are fixed, while query selectivity is kept
constant. Hence, the number of matching records increases proportionally
with the index size, allowing the experiment to capture both encrypted
index lookup and output-sensitive result processing under increasing
data volume. For Scheme~\cite{ref30}, this setting is consistent with
its output-sensitive search complexity $O(n_w^l)$.

\begin{figure}[!t]
    \centering
    \includegraphics[width=\columnwidth]{images/fig_exp2_search.pdf}
    \caption{Encrypted-search latency versus searchable-index size.}
    \label{fig:exp2}
\end{figure}

Fig.~\ref{fig:exp2} shows that search latency increases with index size
for all evaluated schemes, while MA-LB-PQ-VDSE exhibits slower latency
growth. Before search execution, the Authorization Index Manager
validates the current VAP and derives the policy-state-aware token set
$\mathcal{T}_Q$ together with the required authorized shard set
$\mathcal{S}_Q$. AASS then assigns each required shard only to an
eligible FSN maintaining that shard. Consequently, online search is
restricted to policy-authorized shards rather than traversing the
global searchable repository.

For a fixed query and authorization scope, the dominant index-access
cost is approximately
\[
O(|\mathcal{T}_Q|)T_L
\]
while the measured latency additionally grows with the number of matching
entries returned under the constant-selectivity workload. Thus, increasing
the global index size does not proportionally increase the number of
index partitions that must be searched; instead, the main growth arises
from output-sensitive result processing within the selected shards.

Scheme~\cite{ref30} scales with the number of accessible matching
records, Scheme~\cite{ref54} combines encrypted-index lookup with
candidate filtering, while Schemes~\cite{ref35} and~\cite{ref41}
incur their respective query-dependent cryptographic processing.
Overall, the results show that policy-state-aware token lookup and
authorized shard selection limit unnecessary index processing and
improve encrypted-search scalability over large EHR collections.


\subsubsection*{\textbf{Experiment 3: Cross-Domain Search Scalability}}
\label{exp:crossdomain}

This experiment evaluates encrypted-search latency as the number of
participating healthcare domains $d$ increases from 2 to 10. The
proposed framework is compared with Schemes~\cite{ref30},
\cite{ref35}, \cite{ref41}, and \cite{ref54}. Since these baselines do
not provide equivalent authorization-aware cross-domain orchestration,
their domain-local searches are independently executed and the results
aggregated. The query size and per-domain index size are fixed to isolate
cross-domain search overhead.

\begin{figure}[!t]
    \centering
    \includegraphics[width=\columnwidth]{images/fig_exp3_crossdomain.pdf}
    \caption{Cross-domain search latency versus number of participating domains.}
    \label{fig:exp3}
\end{figure}

Fig.~\ref{fig:exp3} shows increasing latency as more domains participate,
with MA-LB-PQ-VDSE exhibiting slower growth. The AIM first validates the
current VAP and derives policy-state-aware tokens and the required
authorized shard set $\mathcal{S}_Q$. AASS then assigns each required
shard to an eligible FSN maintaining that shard, so search execution is
limited to authorization-valid index partitions rather than all
participating domain indexes.

As $d$ increases, the baselines accumulate domain-local search overhead
because their search procedures are independently executed across the
participating domains. In contrast, MA-LB-PQ-VDSE restricts execution to
the shards required by the authorized query and exploits distributed
shard processing through AASS. The results therefore demonstrate better
cross-domain scalability by reducing unnecessary index processing as the
number of participating healthcare domains increases.
\subsubsection*{\textbf{Experiment 4: Fine-Grained Verification Effectiveness}}
\label{exp:verify}

This experiment evaluates the overhead and granularity of verifiable
retrieval as the number of returned entries $r$ increases from 10 to
1000. Controlled tampered or stale entries are also injected to evaluate
whether invalid results can be individually identified without rejecting
otherwise valid results. We report the average verification latency
$T_{\rm avg}=T_{\rm verify}/r$, invalid-result localization rate, and
valid-result retention. The proposed framework is compared with
Schemes~\cite{ref30}, \cite{ref35}, and \cite{ref54}, which provide
explicit retrieval-verification mechanisms.

\begin{figure}[!t]
    \centering
    \includegraphics[width=\columnwidth]{images/fig_exp4_verify.pdf}
    \caption{Fine-grained verification effectiveness:
    (a) average verification latency per returned entry and
    (b) invalid-result localization and valid-result retention.}
    \label{fig:exp4}
\end{figure}

Fig.~\ref{fig:exp4}(a) shows that MA-LB-PQ-VDSE incurs a moderate
per-result verification cost. For each returned entry $I_{i,j}$, the DU
reconstructs
\begin{equation}
L_{i,j}=H(\mathsf{Encode}(I_{i,j}))
\end{equation}
and verifies its Merkle path against $Root_i$. It then checks the latest
blockchain-anchored metadata and recomputes
\begin{equation}
Commit_i^*
=
H(CID_i\parallel Root_i\parallel PID_i
\parallel PV_i\parallel AuthState_i)
\end{equation}
The resulting verification cost is dominated by
$O(r\log t)T_H+O(r)T_{\rm BC}$, where $t$ denotes the number of indexed
keywords in the corresponding record-level Merkle tree.

Fig.~\ref{fig:exp4}(b) demonstrates the benefit of per-entry verification.
Because each returned entry is accompanied by its own authentication
evidence, a modified, stale, or policy-state-inconsistent entry can be
localized and rejected while valid entries remain acceptable. In
contrast, schemes relying primarily on aggregated verification evidence
provide less direct fault localization at individual-result granularity.

Overall, MA-LB-PQ-VDSE introduces additional per-result verification
overhead in exchange for fine-grained integrity and freshness checking.
This enables corrupted or stale results to be isolated without
discarding the entire response, which is particularly useful for
cross-domain IoMT retrieval.



\subsubsection*{\textbf{Experiment 5: Dynamic Index Update}}
\label{exp:update}

This experiment evaluates dynamic index-update latency as the number of
modified keyword--record pairs $k$ increases from $10^2$ to $10^5$.
MA-LB-PQ-VDSE is compared with Scheme~\cite{ref30} (Peony++) and
Scheme~\cite{ref35}, which explicitly support dynamic searchable-index
updates. The total index size is fixed to isolate the effect of increasing
update volume.

\begin{figure}[!t]
    \centering
    \includegraphics[width=\columnwidth]{images/fig_exp5_update.pdf}
    \caption{Dynamic index-update latency versus number of modified
    keyword--record pairs.}
    \label{fig:exp5}
\end{figure}

Fig.~\ref{fig:exp5} shows increasing update latency for all schemes,
while MA-LB-PQ-VDSE maintains lower growth as the update batch increases.
Scheme~\cite{ref30} follows its $O(k)$ update behavior, whereas
Scheme~\cite{ref35} additionally incurs per-keyword hash/PRF processing
and puncturable-PRF operations.

For MA-LB-PQ-VDSE, an update modifies only the affected searchable-index
entries and their Merkle authentication paths. For each changed entry,
the policy-state-aware token and leaf
$L_{i,j}=H(\mathsf{Encode}(I_{i,j}))$ are recomputed, followed by
incremental Merkle maintenance. The dominant cost is therefore
approximately
\[
O(k\log t)T_H
\]
where $t$ denotes the number of indexed keywords in the corresponding
record-level Merkle tree. Unchanged entries and unrelated records require
no reconstruction.

The results demonstrate that localized index evolution and incremental
authenticated-state maintenance enable MA-LB-PQ-VDSE to efficiently
support large dynamic update batches without rebuilding the complete
searchable index.



\subsubsection*{\textbf{Experiment 6: DIAS Synchronization Ablation}}
\label{exp:sync}

This experiment evaluates the effectiveness of Dependency-Aware
Incremental Authorization Synchronization (DIAS) as the fraction of
policies affected by an authority-state update increases from 10\% to
100\%. Three configurations are compared: (i) \emph{Full-State
Synchronization}, which reconstructs and propagates the relevant
authorization/index state to all FSNs; (ii) \emph{Incremental-All},
which updates only affected state but propagates the resulting deltas
to all FSNs; and (iii) \emph{DIAS}, which updates only dependent
policies, records, index entries, and Merkle paths and propagates the
deltas only to FSNs maintaining affected shards. Synchronization
latency and transferred synchronization data are measured.

\begin{figure}[!t]
    \centering
    \includegraphics[width=\columnwidth]
    {images/fig_exp6_sync.pdf}
    \caption{DIAS synchronization under increasing affected-policy ratio.}
    \label{fig:exp6}
\end{figure}

Fig.~\ref{fig:exp6} shows that DIAS achieves the lowest synchronization
overhead, particularly when authority updates affect only a small
fraction of policies. An update by $AA_k$ is restricted through the
dependency chain
\begin{equation}
AA_k
\rightarrow\mathcal{P}_k^{\rm aff}
\rightarrow\mathcal{R}_k^{\rm aff}
\rightarrow\mathcal{S}_k^{\rm aff}
\rightarrow\mathcal{F}_k^{\rm aff}
\end{equation}
so unrelated policies, records, shards, and FSNs require no update.
Incremental-All reduces state-reconstruction overhead but still incurs
unnecessary propagation, whereas Full-State Synchronization processes
substantially more state.

As the affected-policy ratio approaches 100\%, the advantage of DIAS
narrows because a larger portion of the system becomes dependency
relevant. The results therefore confirm that DIAS gains its efficiency
from both dependency-localized state maintenance and selective
propagation, with the largest benefit under sparse authorization
changes typical of multi-authority environments.



\subsubsection*{\textbf{Experiment 7: AASS Search Throughput}}
\label{exp:throughput}

This experiment evaluates AASS under increasing concurrent search
workloads. Four scheduling strategies are compared: (i) \emph{No Load
Balancing}; (ii) \emph{Round Robin}; (iii) \emph{Least Loaded}, which
considers current workload but not shard locality or synchronization
state; and (iv) \emph{AASS}. Concurrent queries increase from 100 to
5000, while the number of FSNs and per-query workload are fixed.

For each query, AASS assigns every required shard
$S\in\mathcal{S}_Q$ only to eligible FSNs maintaining that shard. Among
replicated eligible nodes, it selects the FSN minimizing
\begin{equation}
SC_j=
\lambda_1C_j^{index}
+\lambda_2C_j^{verify}
+\lambda_3C_j^{sync}
+\lambda_4C_j^{queue}
\end{equation}

\begin{figure}[!t]
    \centering
    \includegraphics[width=\columnwidth]{images/fig_exp7_throughput.pdf}
    \caption{Search throughput under increasing concurrent workloads.}
    \label{fig:exp7}
\end{figure}

Fig.~\ref{fig:exp7} shows that AASS sustains substantially higher
throughput than unbalanced scheduling as concurrency increases, reaching
982 queries/s against 543 at 10{,}000 concurrent queries, while Least
Loaded attains the highest raw throughput of the four (1{,}030
queries/s). Without load balancing, requests concentrate on individual
FSNs and cause early saturation. Round Robin improves workload
distribution but ignores shard locality and current node state, while
Least Loaded reduces queue congestion without considering whether the
selected FSN maintains the required shard or requires synchronization.

AASS instead restricts each shard assignment to eligible FSNs and selects
the node with the lowest predicted search cost. By jointly considering
index-processing, verification, synchronization, and queueing overhead,
AASS reduces unnecessary forwarding and avoids overloaded or
synchronization-expensive nodes. The resulting trade is explicit: AASS
accepts a small throughput margin relative to Least Loaded in exchange
for eliminating cross-node forwarding entirely
(Fig.~\ref{fig:exp8}(c)), which Least Loaded cannot do because it
disregards shard locality. Combining locality with workload- and
state-aware scheduling therefore buys correctness of placement at a
bounded throughput cost, rather than dominating on throughput alone.

\subsubsection*{\textbf{Experiment 8: Load-Balancing Effectiveness}}
\label{exp:balance}

This experiment further evaluates AASS by examining workload distribution
across FSNs as concurrent queries increase from 100 to 10,000. The four
strategies in Experiment~\ref{exp:throughput} are compared using:
(i) standard deviation of FSN utilization; (ii) maximum FSN utilization;
and (iii) cross-node forwards. A cross-node forward occurs when a
scheduler assigns a required shard to an FSN that does not maintain it,
requiring redirection to an eligible node.

\begin{figure}[!t]
    \centering
    \includegraphics[width=\columnwidth]{images/fig_exp8_balance.pdf}
    \caption{Load-balancing effectiveness under increasing concurrent
    workloads: (a) utilization standard deviation, (b) maximum FSN
    utilization, and (c) cross-node forwards.}
    \label{fig:exp8}
\end{figure}

Fig.~\ref{fig:exp8}(a) shows that AASS maintains substantially lower
utilization variation than unbalanced scheduling (0.090 against 0.392 at
10{,}000 concurrent queries) and improves as load grows, indicating more
balanced workload distribution across eligible FSNs. Least Loaded attains
the lowest variation of the four (0.019), since it optimises queue
occupancy directly and is unconstrained by shard locality; AASS is
constrained to eligible FSNs and is therefore balancing over a smaller
candidate set at each assignment. Round Robin and Least Loaded improve
upon no load balancing but do not jointly account for required-shard
locality, synchronization state, and predicted processing cost.
Consequently, Fig.~\ref{fig:exp8}(b) shows that AASS also limits
processing hot spots as concurrency increases.

The locality benefit is further demonstrated in
Fig.~\ref{fig:exp8}(c). Conventional strategies may assign work to an
FSN that does not maintain the required shard, causing cross-node
forwarding. AASS instead considers only FSNs satisfying
$S\in\mathcal{S}_j$ for each required shard $S\in\mathcal{S}_Q$ and
selects the eligible node with the lowest predicted cost. It therefore
eliminates unnecessary forwarding -- measuring zero cross-node forwards
at every concurrency level, against approximately 2{,}100 for each of the
three conventional strategies at 10{,}000 concurrent queries -- while
maintaining balanced utilization. Together with
Experiment~\ref{exp:throughput}, these results show that AASS improves
distributed-search scalability through shard-local, state-aware, and
workload-aware scheduling, and that its distinguishing property is the
elimination of misplaced shard assignments rather than a uniform
advantage on throughput or utilization spread.

\section{Conclusion}
\label{sec:conclusion}

This paper presented \textbf{MA-LB-PQ-VDSE}, a post-quantum
authorization-aware framework for secure and scalable cross-domain
encrypted EHR retrieval in IoMT environments. The framework integrates
Version-Bound Authorization Profiles (VAPs), Policy-Driven Searchable
Indexes (PDSIs), authorization-aware search resolution, Adaptive
Authorization-Aware Search Scheduling (AASS), Incremental Authorization
Synchronization (IAS), and blockchain-anchored verification into a
unified retrieval architecture. By resolving authorization before
encrypted matching, MA-LB-PQ-VDSE restricts search execution to
permitted domains and effective candidate sets, while AASS exploits
authorization validity, index locality, synchronization state, and
workload information to improve distributed search execution. IAS
further limits authorization-update overhead through incremental state
maintenance and selective propagation to affected Fog Search Nodes. The security and performance analyses demonstrate that these mechanisms
provide scalable dynamic search while preserving authorization
consistency and verifiability. In particular, the proposed verification
mechanism associates independent authenticated evidence with each
returned ciphertext, enabling a tampered, stale, or authorization-inconsistent result to be individually identified and
rejected rather than invalidating the entire result set. Experimental
results further demonstrate scalable trapdoor generation, encrypted
search, cross-domain retrieval, authenticated dynamic updates, and
authorization synchronization. The AASS experiments also show improved
search throughput and workload distribution under increasing
concurrency, while fine-grained per-result verification provides
stronger validation at a moderate additional verification cost. Future work will investigate
larger-scale heterogeneous IoMT deployments, adaptive optimization of
authorization and scheduling policies, and further reduction of
verification and synchronization overhead for highly dynamic
cross-domain environments.


\section{Acknowledgment}
This work was supported by the Thammasat University Research Fund under Contract TUFT 032/2568.
\bibliographystyle{IEEEtran}


\begin{thebibliography}{99}

\bibitem{ref1}
J. Malamas, V. Kotzanikolaou, P. Dasaklis, and M. Burmester, “A hierarchical multi blockchain for fine grained access to medical data,” \emph{IEEE Access}, vol. 8, pp. 134393–134412.

\bibitem{ref2}
H. Gao, Z. Ma, S. Luo, Y. Xu, and Z. Wu, “BSSPD: a blockchain-based security sharing scheme for personal data with fine-grained access control,” \emph{Wireless Communications and Mobile Computing}, 2021.

\bibitem{ref3}
S. Fugkeaw, “A Lightweight Policy Update Scheme for Outsourced Personal Health Records Sharing,” \emph{IEEE Access}, vol. 9, pp. 54862–54871, 2021.

\bibitem{ref4}
J. Chen, X. Yin, and J. Ning, “A fine-grained and secure health data sharing scheme based on blockchain,” \emph{Transactions on Emerging Telecommunications Technologies}, e4510.

\bibitem{ref5}
J. Yuan, Y. Ma, W. Luo, and G. Han, “B-SSMD: A Fine-Grained Secure Sharing Scheme of Medical Data Based on Blockchain,” \emph{Security and Communication Networks}, 2022.

\bibitem{ref6}
Z. Zhang, J. Wang, Y. Wang, Y. Su, and X. Chen, “Towards efficient verifiable forward secure searchable symmetric encryption,” in \emph{Proc. Eur. Symp. Res. Comput. Secur.}, Sep. 2019, pp. 304–321.

\bibitem{ref7}
Z. Liu, T. Li, P. Li, C. Jia, and J. Li, “Verifiable searchable encryption with aggregate keys for data sharing system,” \emph{Future Gener. Comput. Syst.}, vol. 78, pp. 778–788, Jan. 2018. [Online]. Available: http://dx.doi.org/10.1016/j.future.2017.02.024

\bibitem{ref8}
Z. Shi, X. Fu, X. Li, and K. Zhu, “ESVSSE: Enabling efficient, secure, verifiable searchable symmetric encryption,” \emph{IEEE Trans. Knowl. Data Eng.}, vol. 34, no. 7, pp. 3241–3254, Jul. 2022, doi: 10.1109/TKDE.2020.3025348.

\bibitem{ref9}
A. Wu, A. Yang, W. Luo, and W. Jinghang, “Enabling traceable and verifiable multi-user forward secure searchable encryption in hybrid cloud,” \emph{IEEE Trans. Cloud Comput.}, early access, Apr. 26, 2022.

\bibitem{ref10}
Q. Gan, J. K. Liu, X. Wang, \emph{et al.}, “Verifiable searchable symmetric encryption for conjunctive keyword queries in cloud storage,” \emph{Front. Comput. Sci.}, 16, 166820 (2022).

\bibitem{ref11}
S. Hu, C. Cai, Q. Wang, C. Wang, X. Luo, and K. Ren, “Searching an encrypted cloud meets blockchain: A decentralized, reliable and fair realization,” in \emph{Proc. IEEE Conf. Comput. Commun. (INFOCOM)}, Apr. 2018, pp. 792–800.

\bibitem{ref12}
T. Wang \emph{et al.}, “An efficient verifiable searchable encryption scheme with aggregating authorization for blockchain-enabled IoT,” \emph{IEEE Internet Things J.}, vol. 9, no. 20, pp. 20666–20680, Oct. 2022.

\bibitem{ref13}
R. Du, N. Liu, M. Li, and J. Tian, “Block verifiable dynamic searchable encryption using redactable blockchain,” \emph{Journal of Information Security and Applications}, vol. 75, 2023.

\bibitem{ref14}
B. Chen, T. Xiang, D. He, H. Li and K.-K. R. Choo, “BPVSE: Publicly Verifiable Searchable Encryption for Cloud-Assisted Electronic Health Records,” \emph{IEEE Transactions on Information Forensics and Security}, vol. 18, pp. 3171–3184, 2023.

\bibitem{ref15}
C. Huang, D. Liu, A. Yang, R. Lu and X. Shen, “Multi-client Secure and Efficient DPF-based Keyword Search for Cloud Storage,” \emph{IEEE Transactions on Dependable and Secure Computing}.

\bibitem{ref16}
L. Xu, X. Yuan, R. Steinfeld, C. Wang, and C. Xu, “Multi-writer searchable encryption: An lwe-based realization and implementation,” in \emph{Proc. of ACM CCS}, 2019, pp. 122–133.

\bibitem{ref17}
J. Gharehchamani, Y. Wang, D. Papadopoulos, M. Zhang, and R. Jalili, “Multi-user dynamic searchable symmetric encryption with corrupted participants,” \emph{IEEE Trans. Dependable Secure Comput.}, 2021.

\bibitem{ref18}
X. Zhang, C. Huang, Y. Su, J. Qin, and X. Shen, “Divertible searchable symmetric encryption for secure cloud storage,” in \emph{Proc. of IEEE GLOBECOM}, 2022, pp. 3785–3790.

\bibitem{ref19}
X. Yang, G. Chen, M. Wang, T. Li and C. Wang, “Multi-Keyword Certificateless Searchable Public Key Authenticated Encryption Scheme Based on Blockchain,” \emph{IEEE Access}, vol. 8, pp. 158765–158777, 2020.

\bibitem{ref20}
J. Lapmoon and S. Fugkeaw, “A Verifiable and Secure Industrial IoT Data Deduplication Scheme With Real-Time Data Integrity Checking in Fog-Assisted Cloud Environments,” \emph{IEEE Access}, vol. 13, pp. 11969–11988, 2025.

\bibitem{ref21}
S. Fugkeaw, R. Prasad Gupta and K. Worapaluk, “Secure and Fine-Grained Access Control With Optimized Revocation for Outsourced IoT EHRs With Adaptive Load-Sharing in Fog-Assisted Cloud Environment,” \emph{IEEE Access}, vol. 12, pp. 82753–82768, 2024.

\bibitem{ref22}
F. Li \emph{et al.}, “Towards Efficient Verifiable Boolean Search Over Encrypted Cloud Data,” \emph{IEEE Transactions on Cloud Computing}, vol. 11, no. 1, pp. 839–853, 1 Jan.–Mar. 2023.

\bibitem{ref23}
Q. Wang, X. Zhang, J. Qin, J. Ma, and X. Huang, “A Verifiable Symmetric Searchable Encryption Scheme Based on the AVL Tree,” \emph{The Computer Journal}, vol. 66, no. 1, Jan. 2023, pp. 174–183, 

\bibitem{ref24}
L. Ji, J. Li, Y. Zhang and Y. Lu, “Verifiable Searchable Symmetric Encryption Over Additive Homomorphism,” \emph{IEEE Transactions on Information Forensics and Security}, vol. 20, pp. 1320–1332, 2025.

\bibitem{ref25}
S. Fugkeaw, L. Hak and T. Theeramunkong, “Achieving Secure, Verifiable, and Efficient Boolean Keyword Searchable Encryption for Cloud Data Warehouse,” \emph{IEEE Access}, vol. 12, pp. 49848–49864, 2024.

\bibitem{ref26}
L. Xu, S. Sun, X. Yuan, J. K. Liu, C. Zuo and C. Xu, “Enabling Authorized Encrypted Search for Multi-Authority Medical Databases,” \emph{IEEE Transactions on Emerging Topics in Computing}, vol. 9, no. 1, pp. 534–546, 1 Jan.–Mar. 2021.

\bibitem{ref27}
J. Bethencourt, A. Sahai, and B. Waters, “Ciphertext-policy attribute-based encryption,” in \emph{Proc. IEEE Symp. Security and Privacy}, 2007, pp. 321–334.

\bibitem{ref28}
J. Yu, S. Liu, M. Xu, H. Guo, F. Zhong and W. Cheng, “An Efficient Revocable and Searchable MA-ABE Scheme With Blockchain Assistance for C-IoT,” \emph{IEEE Internet of Things Journal}, vol. 10, no. 3, pp. 2754–2766, 1 Feb. 2023.

\bibitem{ref29}
H. Dou \emph{et al.}, “Dynamic Searchable Symmetric Encryption With Strong Security and Robustness,” \emph{IEEE Transactions on Information Forensics and Security}, vol. 19, pp. 2370–2384, 2024.

\bibitem{ref30}
Y. Ge, Y. Gao, J. Ning, J. Ma and X. Chen, “Verifiable Multilevel Dynamic Searchable Encryption With Forward and Backward Privacy in Cloud-Assisted IoT,” \emph{IEEE Internet of Things Journal}, vol. 11, no. 24, pp. 40861–40874, 15 Dec. 2024.

\bibitem{ref31}
D. Li, X. Zhao, H. Li and K. Fan, “Volume-Hiding Multidimensional Verifiable Dynamic Searchable Symmetric Encryption Scheme for Cloud Computing,” \emph{IEEE Internet of Things Journal}, vol. 11, no. 23, pp. 37437–37451, 1 Dec. 2024.

\bibitem{ref32}
X. Zhu, J. Zhou, Y. Dai, P. Shen, S. Kasra Kermanshahi and J. Hu, “A Verifiable and Efficient Symmetric Searchable Encryption Scheme for Dynamic Dataset With Forward and Backward Privacy,” \emph{IEEE Transactions on Dependable and Secure Computing}, vol. 22, no. 3, pp. 2741–2755, May–Jun. 2025.

\bibitem{ref33}
M. Li, C. Jia, R. Du and W. Shao, “Forward and Backward Secure Searchable Encryption Scheme Supporting Conjunctive Queries Over Bipartite Graphs,” \emph{IEEE Transactions on Cloud Computing}, vol. 11, no. 1, pp. 1091–1102, 1 Jan.–Mar. 2023.

\bibitem{ref34}
M. Li, C. Jia, R. Du, W. Shao and G. Ha, “DSE-RB: A Privacy-Preserving Dynamic Searchable Encryption Framework on Redactable Blockchain,” \emph{IEEE Transactions on Cloud Computing}, vol. 11, no. 3, pp. 2856–2872, Jul.–Sep. 2023.

\bibitem{ref35}
C. Guo, W. Li, X. Tang, K.-K. R. Choo and Y. Liu, “Forward Private Verifiable Dynamic Searchable Symmetric Encryption With Efficient Conjunctive Query,” \emph{IEEE Transactions on Dependable and Secure Computing}, vol. 21, no. 2, pp. 746–763, Mar.–Apr. 2024.

\bibitem{ref36}
Q. Jiang, X. Yang, S. Chen, S. Qi, F. Song and Z. Fu, “XB-Muse: Practical Multiuser Dynamic Searchable Symmetric Encryption for Adaptive Revocation,” \emph{IEEE Internet of Things Journal}, vol. 12, no. 14, pp. 26486–26499, 15 Jul. 2025, doi: 10.1109/JIOT.2025.3561287.

\bibitem{ref37}
X. Wang, J. Ma, X. Liu, Y. Miao, Y. Liu and R. H. Deng, “Forward/Backward and Content Private DSSE for Spatial Keyword Queries,” \emph{IEEE Transactions on Dependable and Secure Computing}, vol. 20, no. 4, pp. 3358–3370, Jul.–Aug. 2023.

\bibitem{ref38}
X. Zhang, C. Huang, Y. Su and J. Qin, “Secure, Dynamic, and Efficient Keyword Search With Flexible Merging for Cloud Storage,” \emph{IEEE Transactions on Services Computing}, vol. 17, no. 5, pp. 2822–2835, Sep.–Oct. 2024.

\bibitem{ref39}
D. X. Song, D. Wagner, and A. Perrig, “Practical techniques for searches on encrypted data,” in \emph{Proc. IEEE Symp. Security and Privacy}, 2000, pp. 44–55.

\bibitem{ref40}
R. Curtmola, J. Garay, S. Kamara, and R. Ostrovsky, “Searchable symmetric encryption: Improved definitions and efficient constructions,” in \emph{Proc. ACM CCS}, 2006, pp. 79–88.

\bibitem{ref41}
C. Thingom et al., ”Secure and Privacy-Preserving Post-Quantum
Attribute-Based Searchable Encryption for Edge-Driven Transportation
Systems,” IEEE Transactions on Consumer Electronics, vol. 72, no. 1,
pp. 1865–1875, Feb. 2026, doi: 10.1109/TCE.2025.3632071.

\bibitem{ref42}
Y. Yang, J. Ma, X. Liu, and R. H. Deng, “Revocable dynamic searchable encryption with forward and backward privacy,” \emph{IEEE Transactions on Dependable and Secure Computing}, early access, 2025.

\bibitem{ref43}
B. Chen, T. Xiang, D. He, H. Li, and K.-K. R. Choo, “BPVSE: Publicly Verifiable Searchable Encryption for Cloud-Assisted Electronic Health Records,” \emph{IEEE Transactions on Information Forensics and Security}, vol. 18, pp. 3171–3184, 2023.

\bibitem{ref44}
S. Fugkeaw, K. Tangtanawirut, P. Rattanasrisuk and A. Changtor, "MK-WISE: Secure and Efficient Multi-Keyword Wildcard ABSE with Keyword-Level Revocation for Device–Edge–Cloud EHRs Data Sharing," \emph {IEEE Transactions on Network and Service Management}.

\bibitem{ref45}
J. Wu, K. Zhang, J. Ning, H. Chen, L. Zhang and Z. Ying, "Multi-Writer/Reader Forward and Backward Private DSSE With Bilateral Selection," in IEEE Transactions on Dependable and Secure Computing, vol. 23, no. 2, pp. 2444-2457, March-April 2026, doi: 10.1109/TDSC.2025.3627926.

\bibitem{ref46}
L. Chen, Y. Wu, H. Zhang, J. Weng and Z. Liu, "Verifiable Multi-User Dynamic Searchable Symmetric Encryption With Forward and Backward Privacy Feasible for Cloud Storage," in IEEE Transactions on Information Forensics and Security, vol. 21, pp. 3797-3811, 2026, doi: 10.1109/TIFS.2026.3681434.

\bibitem{ref47}
H. Gao, H. Huang, L. Xue, F. Xiao, and Q. Li, “Blockchain-
enabled fine-grained searchable encryption with cloud–edge com-
puting for electronic health records sharing,” IEEE Internet of
Things Journal, vol. 10, no. 20, pp. 18414–18425, Oct. 2023, doi:
10.1109/JIOT.2023.3279893.

\bibitem{ref48}
M. Perera and S. Fugkeaw, "Enabling Fault-Tolerant and Load-Balanced Multi-Query Searchable Encryption With Multi-Level Revocation in Edge-Fog-Cloud for Electronic Health Records," in IEEE Open Journal of the Communications Society, vol. 7, pp. 6964-6983, 2026, doi: 10.1109/OJCOMS.2026.3706304. 

\bibitem{ref49}
Y. Miao et al., “Time-controllable keyword search scheme with
efficient revocation in mobile e-health cloud,” IEEE Transactions on
Mobile Computing, vol. 23, no. 5, pp. 3650–3665, May 2024.

\bibitem{ref50}
L. Hak and S. Fugkeaw, "SSL-XIoMT: Secure, Scalable, and Lightweight Cross-Domain IoMT Sharing With SSI and ZKP Authentication," in IEEE Open Journal of the Computer Society, vol. 6, pp. 714-725, 2025, doi: 10.1109/OJCS.2025.3570087.

\bibitem{ref51}
B. D. Deebak and S. O. Hwang, "Blockchain Authorized Privacy-Preserving Framework Using Edge Computing for the Cloud-Assisted Internet of Medical Things," in IEEE Transactions on Cloud Computing, vol. 14, no. 2, pp. 445-462, April-June 2026, doi: 10.1109/TCC.2026.3656828. 

\bibitem{ref52}
E.-S. Zhuang, C.-Y. Lin, C.-I. Fan, A. Karati, and D. Das, “Multi-
authority attribute-based multi-keyword searchable encryption with dy-
namic membership from lattices,” in Proc. IEEE Conf. Dependable
and Secure Computing (DSC), Taipei, Taiwan, 2025, pp. 1–6, doi:
10.1109/DSC65356.2025.11260864.

\bibitem{ref53}
L. Chen, Y. Liu, H. Zhang and J. Weng, ”Multi-Authority Attribute-
Based Searchable Encryption With Attribute Revocation for Internet of
Vehicles,” IEEE Transactions on Vehicular Technology, early access,
2025, doi: 10.1109/TVT.2025.3629863.

\bibitem{ref54}
M. Perera and S. Fugkeaw,
``LV-PQ-ABSE: A Lightweight Verifiable Post-Quantum Attribute-Based
Searchable Encryption Scheme with Hybrid Indexing and
Provenance-Aware Verification for IoT-Based EHRs,''
\textit{IEEE Internet of Things Journal}, early access, 2026,
doi: 10.1109/JIOT.2026.3695855.

\bibitem{ref55}
Y. Cao et al.,
``Enabling Puncturable Encrypted Search Over Lattice for
Privacy-Preserving in Mobile Cloud,''
\textit{IEEE Transactions on Mobile Computing}, early access, 2026,
doi: 10.1109/TMC.2026.3682118.

\bibitem{ref56}
A. Johnson, L. Bulgarelli, T. Pollard, B. Gow, B. Moody,
S. Horng, L. A. Celi, and R. Mark,
``MIMIC-IV (version 3.1),''
PhysioNet, 2024, doi: 10.13026/kpb9-mt58.

\end{thebibliography}



\end{document}
